% =============================================================================
% The Bridge Certificate -- a line-level audit instrument for inter-theoretic
% transfer arguments, with four applications to the abc literature.
% Draft v0.2 (English) -- leading language version.
% Project: abc_forensik/paper_bridge_certificate  (.RESEARCH/.LAB/.HCT)
% Part 3 of the abc Gap Series (numbering per the series register, decision E2=a
% of 2026-08-13). This is NOT a proof paper and NOT an error claim.
% See the scope box after the abstract.
%
% v0.2 (2026-08-16): adopts the hardened sixth contract component of the now
% published methods paper (Part 2, v0.2.4). Two verdicts of v0.1 are recoded --
% Case A obligation 6 (discharged -> does not deliver) and Case C obligation 6
% (discharged -> partly) -- and the recoding is declared openly in the
% unnumbered section "Changes in version 0.2" after the table of contents.
% Further: R6a/R6b split, version label and row-namespace note for the
% instrument, R8 clarified as a triggered extension (E8), Case C relabelled
% "comparative discrimination case", LANA mapping booked primarily at
% obligation 6, Part 2 cited as published. v0.1 remains unchanged on Zenodo.
% =============================================================================

\documentclass[11pt,a4paper]{article}

\usepackage[utf8]{inputenc}
\usepackage{lmodern}   % scalable fonts -- required for microtype font expansion
\usepackage{cmap}
\usepackage[T1]{fontenc}
\usepackage[english]{babel}
\usepackage[a4paper,margin=2.7cm]{geometry}
\usepackage{amsmath}
\usepackage{amsthm}
\usepackage{amssymb}
\usepackage{array}
\usepackage{booktabs}
\usepackage{microtype}
\usepackage{xurl}
\usepackage[colorlinks=true,linkcolor=black,citecolor=black,urlcolor=blue]{hyperref}

\input{glyphtounicode}
\pdfgentounicode=1

\hyphenation{Teich-mueller inter-uni-ver-sal in-de-ter-mi-na-cy
  in-de-ter-mi-na-cies multi-ra-di-al pro-ces-sion nor-mal-i-za-tion
  arith-me-ti-coid arith-me-ti-coids}

\emergencystretch=1.5em

\newcolumntype{P}[1]{>{\raggedright\arraybackslash}p{#1}}

\hypersetup{
  pdfencoding=unicode,
  psdextra,
  pdftitle={The Bridge Certificate},
  pdfauthor={Lukas Geiger},
  pdfsubject={A line-level audit instrument for inter-theoretic transfer arguments},
  }

% Language-specific destination prefix, so that the bilingual combined PDF keeps
% the anchors of the two language versions apart.
\makeatletter
\renewcommand*{\HyperDestNameFilter}[1]{bc-en-#1}
\makeatother

% --- environments -------------------------------------------------------------
\theoremstyle{definition}
\newtheorem{obligation}{Obligation}
\newtheorem*{workinglemma}{Working lemma (bridge contract)}
\newtheorem*{recipe}{Application recipe}

% --- notation -----------------------------------------------------------------
\newcommand{\Brxi}{\mathrm{Br}_{\xi}}
\newcommand{\LogVol}{\operatorname{LogVol}}
\newcommand{\Vol}{\operatorname{Vol}}
\newcommand{\ordd}{\operatorname{ord}}
\newcommand{\Thet}{\widetilde{\Theta}}

% Every source citation is anchored to a named extract, because the extracts
% used here carry independent line numbering; a bare line number would be
% ambiguous.
\newcommand{\anchor}[2]{(#1, extract l.~#2)}

\title{The Bridge Certificate:\\
A Line-Level Audit Instrument for Inter-Theoretic\\
Transfer Arguments, with Four Applications\\
to the $abc$ Literature}
\author{Lukas Geiger\\[2pt]
\small Independent Researcher, Bernau, Germany\\
\small ORCID: 0009-0005-7296-1534}
\date{October 2026 --- Version 0.3 (corrective draft)}

\begin{document}
\maketitle

\begin{abstract}
\noindent
A recurring shape of argument in arithmetic geometry moves a numerical readout
out of one construction into a container it shares with another, and then reads
an inequality off the result. Disputes about such steps have proved unusually resistant to the ordinary
instruments of mathematical review: the parties agree on almost every formula and
disagree about whether a particular passage \emph{is} an argument. This paper
proposes an instrument for exactly that situation, and reports what it returned
in four applications. The instrument --- a \emph{bridge certificate} --- fixes
seven obligations of the proposed certificate for passages from a qualitative
linkage to a numerical membership statement, arranges them as a nine-row ledger
to be filled line by line against the wording of a primary text, and constrains
how findings may be recorded: verdicts are about certificate form, never about
mathematical truth; every finding carries a verbatim quotation with a line anchor;
gaps are typed rather than graded; every appeal to standard literature is entered
in an import ledger; and every claim of the form ``this is not in the text'' is
reported together with the search terms used. Four applications are given, and
they do four different things. Against the proof of one corollary of
inter-universal Teichmüller theory the ledger \emph{locates} an unmet
obligation: the decisive membership is stated to follow ``formally'' from a pair
of declared properties rather than being derived. Against the follow-up paper
that the same text names as its quantification site, the ledger \emph{closes off}
that site: the computation there is explicit and, at every node it names,
self-contained --- this paper did not recompute it --- but it consumes the disputed
inequality as an input rather than producing it. Against an independent
alternative program in a different theory the ledger serves as a
\emph{comparative discrimination case} and returns a differently shaped gap at a
different address. Against two
downstream applications it records \emph{inheritance}: the disputed step is
imported in a single sentence and never discussed. In the same period an
international formalization project, which does not cite this work, isolated a
compatibility statement at the same passage; we map its
objects onto the certificate's obligations by their proposed roles. This is
not a proved implication from the report's compatibility equation to the
certificate's typed order condition. Nothing in this paper is a claim that any
argument examined is wrong, and nothing in it is a statement about the $abc$
conjecture.
\end{abstract}

% -----------------------------------------------------------------------------
% SCOPE AND CLAIM BOUNDARY
% -----------------------------------------------------------------------------
\begin{center}
\fbox{%
\parbox{\dimexpr\textwidth-2\fboxsep-2\fboxrule\relax}{%
\medskip
\textbf{Scope and claim boundary.}\par\medskip
\begin{itemize}\setlength{\itemsep}{2pt}
\item This paper proves \emph{nothing} about the $abc$ conjecture. It contains no
  proof and no disproof of it, and no statement about whether it is provable by
  the methods discussed.
\item It contains \emph{no error claim}. A ledger row returned as ``not
  delivered'' records that a stated obligation is not discharged \emph{in the
  wording examined}. That is a statement about the form of a certificate, not
  about the truth of a theorem. ``Not delivered'' is not ``refuted''.
\item It \emph{verifies no bound}. No proof line was followed step by step, no
  estimate was recomputed, no constant was checked. Where a computation is
  reported as carried out in a source, that is a finding \emph{that} it is
  carried out and \emph{where}, not a confirmation that it is correct.
\item It takes no side in any existing dispute about the arguments examined. The
  known objections enter this paper only as the definition of a negative control,
  never as evidence.
\item \textbf{No authority argument.} Neither an author's own assertion that a
  step is sufficient, nor a critic's assertion that it is not, counts as evidence
  here. Self-assessments found in the examined texts are objects of examination,
  not support.
\item \textbf{Textual basis} are \texttt{pdftotext} extracts of the source PDFs.
  All line references (``extract l.~NNNN'') are finding aids into those extracts
  and are \emph{not} pagination of the published documents. Unlike in a purely
  prose-level audit, one verdict of Section~\ref{sec:caseA} does depend on
  symbols lost in extraction; the reconstruction is disclosed in
  Section~\ref{sec:reconstruction-A} and the rule governing such cases is part of
  the instrument (Section~\ref{sec:reconstruction}).
\item Bibliographic entries are provisional and carry their provenance; three of
  them are deliberately reduced to what could be verified. See
  Section~\ref{sec:limitations}, item~\ref{lim:biblio}.
\end{itemize}
\medskip
}}
\end{center}

\clearpage
\tableofcontents
\clearpage

% =============================================================================
% CHANGES IN VERSION 0.2 -- unnumbered on purpose: it must not shift the
% section numbering of the published version 0.1.
% =============================================================================
\phantomsection
\section*{Changes in version 0.3}
\addcontentsline{toc}{section}{Changes in version 0.3}
\label{sec:changes03}

This corrective version distinguishes the prospective recipe from the partial
historical applications, corrects the Case-C synopsis reference to ATS IV v2
\S\,6.10/Theorem~6.10.1, and restores the LANA equation's notation
$\eta^q=\eta_S^{\mathrm{anab}}$ (9-1).
The LANA crosswalk is a proposed role correspondence: neither a volume quotient
nor equality of abstract maps by itself proves the certificate's typed
half-line condition. The required object, evaluation and bound remain open.
The seven obligations define the proposed form contract; their universal
necessity for every possible mathematical proof is not established.
Case D concerns the Zhou downstream imports, not a completed C7 positive
control or a demonstrated dependency cycle. Case C does not prove a
mathematical contradiction in ATS.
The previously prepared numbered table captions and two synopsis tables
are retained. English and German texts incorporate the same corrections.
No new mathematical result or independent method validation is claimed.

\medskip

\section*{Changes in version 0.2}
\addcontentsline{toc}{section}{Changes in version 0.2}
\label{sec:changes02}

Version 0.1 of this paper appeared before the methods paper of the
series~\cite{GeigerBCM2026} was finished. That paper hardened the sixth contract
component: a common numerical evaluation in the container no longer discharges
it by itself; the compatibility identities the claim needs --- in particular
$\nu_C \circ i_q = \nu_q$ on the stated source domain --- are mandatory, and a
stipulated invariance is not a proof of one. Version 0.2 adopts that reading,
and adopting it changes two verdicts that version 0.1 recorded. A published
coding should not be revised silently, so both are named here.

\begin{itemize}
\item \textbf{Case~A, obligation~6: discharged $\rightarrow$ does not deliver.}
  Version 0.1 recorded obligation~6 as met ``via its first alternative'',
  because a numerical evaluation exists in the container, and stated that ``the
  gap sits at obligation~7, not at obligation~6''. Under the compatibility
  reading the same evidence yields the opposite coding, and the evidence is the
  paper's own: the ``searched, not found'' block of
  Section~\ref{sec:caseA} already recorded that the identity
  $\nu_C \circ i_q = \nu_q$ across the container identification is not in the
  text, and that the one compatibility lemma which does exist lies on another
  axis. No source was re-examined for this change and no new evidence was
  gathered; what changed is the rule the existing evidence is read against.
\item \textbf{Case~C, obligation~6: discharged $\rightarrow$ partly.} Version
  0.1 recorded obligation~6 as met, again via the first alternative. The
  evaluation there is genuinely numerical, but the compatibility holds only
  between the sums and is place-wise disclaimed by the source itself --- in a
  passage version 0.1 already quoted, one paragraph further on.
\end{itemize}

Three consequences are carried through the text rather than confined to this
note: obligation~6 receives a ledger row of its own (the R6a\,/\,R6b split of
Section~\ref{sec:splits}); the application recipe of Section~\ref{sec:recipe} is
corrected at step~4, so that a re-use of it does not reproduce the superseded
reading; and the summary statements which flanked the Case~A finding with a
discharged obligation~6 --- in the results table, in the role table of
Section~\ref{sec:convergence} and in the conclusion --- are brought into line.

\paragraph{What did not change.} The central finding is untouched: obligation~7
is not discharged at the substep examined, and the membership is stated to
follow ``formally'' from a property pair. The abstract of version 0.1 did not
carry the obligation-6 claim and is unchanged in substance. Cases~B and~D keep
their verdicts, and no case was re-audited.

\paragraph{The remaining changes are terminological}, and align this paper with
the finished methods paper. The instrument is given a version label and its
unprefixed row labels are scoped (Section~\ref{sec:ledger}). Row R8 is
identified as a triggered extension rather than a core row, which corrects the
denominator without changing its coding. And Case~C is renamed from ``comparison
control'' to \emph{comparative discrimination case}, the term Part~2 uses --- a
change of label only, since this paper already recorded that a genuine positive
control was not among the cases examined.

\clearpage

% =============================================================================
\section{Introduction}
\label{sec:intro}
% =============================================================================

\subsection{A disagreement that ordinary review does not resolve}

There is a class of mathematical disagreement in which the usual repair mechanism
fails. It is not that one side has made an arithmetic slip and the other has
found it. It is not that a lemma is false and a counterexample settles the
matter. It is that a passage of text is read by its author as an argument and by
a reader as a restatement, and that both readings survive contact with the
formulas, because the formulas are not in dispute.

The passage at the center of this paper has that character. In the third volume
of inter-universal Teichmüller theory, the proof of a corollary reaches a point
at which two quantities --- a numerical readout attached to one object, and a
half-line of real numbers determined by another --- are brought into a single
evaluation, and an inclusion of the first in the second is then asserted. Around
that inclusion an entire arithmetic consequence is organized. Whether the passage
establishes the inclusion or presupposes it has been argued about for years
without convergence.

Ordinary review prose is a poor instrument here, for a structural reason. Prose
review answers the question ``is this correct?''. When the disagreement is about
whether a given sequence of sentences constitutes a derivation at all, that
question has no stable meaning: each side answers it about a different object.
Worse, prose review permits --- indeed invites --- a move that makes the
disagreement unresolvable in principle, namely the appeal to authority. ``The
author says the construction suffices'' and ``the critics say it does not'' are
both statements about people. Once they enter, the text has stopped being the
subject.

\subsection{What an audit instrument does instead}

The alternative pursued here is to stop asking whether the argument is correct
and start asking a question that a text can answer on its own: \emph{which
obligations does a passage of this kind have to discharge, and which of them does
this particular wording discharge?}

That reformulation buys four things.

\begin{enumerate}
\item \textbf{It is decidable from the text.} Whether a document contains a lemma
  asserting a given implication is a question about the document. Two careful
  readers can disagree about whether a proof is convincing; they will rarely
  disagree about whether a sentence is present.
\item \textbf{It localizes.} A verdict is attached to a numbered obligation and,
  through it, to a line range of a specific text. ``The argument has a gap''
  becomes ``obligation~7 is not discharged at this passage, and here is the
  sentence that stands in its place''.
\item \textbf{It is symmetric.} The same instrument, unchanged, is applied to a
  proposal and to its competitor. A method that returns ``fails'' for one text
  and has no procedure for returning anything about another is not an instrument
  but a verdict looking for a subject.
\item \textbf{It bounds its own claims.} Because the question is about form, the
  answer cannot silently grow into a claim about truth. An unmet obligation is a
  missing certificate component, not a refutation --- and the instrument is
  built so that the distinction cannot be lost in the reporting.
\end{enumerate}

The price is equally clear and is stated here once, because everything in this
paper inherits it: an instrument of this kind can never establish that an
argument is sound. Every obligation may be discharged and the argument may still
be wrong for a reason the instrument does not look for. It can only establish
that a specific, named requirement is or is not met by a specific, quoted
wording.

\subsection{The instrument, in one paragraph}

A \emph{bridge certificate}, written $\Brxi$, is a list of seven obligations
attached to an argumentative step of a particular shape: a step that carries a
numerical readout out of one construction, into a container shared with a second
construction, and then reads an order relation off the result. The obligations
demand, in outline, that both sides carry explicit numerical readouts, that the
shared container come with typed embeddings, that the identification between the
sides be a permitted identification rather than a coincidence of names, that the
indeterminacy budget be split into absorbed and unabsorbed parts, that a numerical
evaluation exist in the container, and --- the last and decisive one --- that the
order relation be \emph{derived} from the permitted identification rather than
asserted alongside it. The obligations are operationalized as a ledger of nine
rows, each a question with a recorded expectation and a recorded risk, to be
filled against a primary text. Around the ledger sit five disciplines that make
its findings auditable: a verdict vocabulary that separates form from truth, a
typology that classifies gaps instead of grading them, an import ledger for every
appeal to standard literature, a rule for reporting what was searched for and not
found, and a rule for disclosing reconstructions where the textual basis is
imperfect. Section~\ref{sec:instrument} states all of this in full.

\subsection{Sources, labels and the anchor convention}
\label{sec:sources}

Seven documents are used, in the form of text extracts of their PDFs: six carry the
four applications, and the seventh is the interim report examined in
Section~\ref{sec:convergence}. The short labels of the left column are used
throughout; every source citation in this paper carries one, because the extracts
have independent line numbering and a bare line number would be ambiguous (see Table~\ref{tab:primary_sources}).

\begin{table}[htbp]
\centering
\small
\renewcommand{\arraystretch}{1.12}
\begin{tabular}{@{}lll@{}}
\toprule
Label & Source & Role in this paper \\
\midrule
IUT-III & \cite{Mochizuki2021c} & Case A: the corollary and its proof \\
IUT-IV & \cite{Mochizuki2021d} & Case B: the declared quantification site \\
ATS-III & \cite{JoshiATSIII} & Case C: construction and volume theory \\
ATS-IV & \cite{JoshiATSIV} & Case C: the $abc$ conclusion and its bridge \\
Zhou-I & \cite{ZhouI} & Case D: first downstream application \\
Zhou-II & \cite{ZhouII} & Case D: second downstream application \\
LANA & \cite{LANA2026} & Section~\ref{sec:convergence}: independent formalization \\
\bottomrule
\end{tabular}
\caption{Primary source corpus and analytical roles in the four case audits.}
\label{tab:primary_sources}
\end{table}

Quotations are verbatim from the extracts; mathematical symbols have been
typeset, ellipses are marked, and every reconstruction of a symbol lost in
extraction is disclosed where it occurs. Emphasis inside quotations is ours
unless stated otherwise.

\paragraph{Language of quotations.} A German version of this paper exists, and it
keeps every quotation in the English original. That is not an omission of the
translation but a consequence of the subject matter: here the wording \emph{is}
the evidence. Whether a text writes ``follows formally'' or proves a theorem is
precisely the question at issue, and translating such a passage would replace the
finding by its rendering. Only the surrounding assessment is translated.

\subsection{What follows}

Section~\ref{sec:instrument} states the instrument in full --- all seven
obligations and all nine rows --- and Section~\ref{sec:recipe} gives the
application recipe, so that it can be executed against a transfer argument other
than the four examined here.
Sections~\ref{sec:caseA}--\ref{sec:caseD} report the four applications.
Section~\ref{sec:convergence} states what they jointly show --- and, at equal
length, what they do not. Section~\ref{sec:limitations} lists the limitations,
Section~\ref{sec:conclusion} concludes.

\paragraph{Place in the series.} This is Part~3 of the \emph{abc} Gap Series.
Part~1 is the published forensic study of the same passage~\cite{GeigerGap2026},
on which Section~\ref{sec:caseA} builds and from which it is delimited there. The
certificate stated here is one instance of a broader gate methodology --- the
instance for numerical membership --- and that broader methodology has since
appeared as Part~2 of the series~\cite{GeigerBCM2026}. Version 0.1 of this paper
assumed nothing from it, because it did not yet exist. This version takes one
thing from it and says where: the hardened reading of obligation~6
(Section~\ref{sec:obligations}), which is why two verdicts of the first version
are recoded here --- see ``Changes in version 0.2'' above. What this paper
contributes is the instance and the four applications reported below,
and those applications fill the ledger only in part: Case~A reports four of its
nine rows, as Sections~\ref{sec:caseA} and~\ref{sec:convergence} state.

% =============================================================================
\section{The bridge certificate}
\label{sec:instrument}
% =============================================================================

\subsection{The question}

The instrument addresses steps of the following shape. Two constructions produce
two objects. The first carries a numerical readout $\nu_q$ --- in the cases below,
a normalized logarithmic volume. The second determines a target set, in the cases
below a closed real half-line $H_\Theta$. The two objects are then exhibited
inside a common container, related by data that is variously described as a
linkage, an identification, a transport or a correspondence, and an assertion of
the form
\[
  \nu_q(\text{first object}) \in H_\Theta
\]
is made. From that membership an inequality follows immediately, and the
inequality is the point of the whole construction.

The question the certificate asks is:

\begin{quote}
When may a qualitative package --- two objects, shared labels, an algorithmic
transport, a common container --- be read as a numerical membership statement in
a concrete real half-line?
\end{quote}

This is deliberately narrower than ``is the argument correct?'' and deliberately
wider than ``is this formula right?''. It is the question that the disagreement is
actually about.

\subsection{The seven obligations}
\label{sec:obligations}

In the proposed specification, an admissible certificate $\Brxi$ contains at
least the following. These are obligations of this chosen form contract, not
proved necessary conditions for every possible mathematical proof.

\begin{obligation}[$q$-side object with readout]
An object $\mathrm{Obj}_q$ together with a concrete log-volumetric readout
$\nu_q$.
\end{obligation}

\begin{obligation}[$\Theta$-side object with target]
An object $\mathrm{Obj}_\Theta$ together with the target half-line
$H_\Theta = \mathbb{R}_{\le -|\log(\Theta)|}$.
\end{obligation}

\begin{obligation}[Common container with typed embeddings]
A container $C$ together with explicitly typed embeddings $i_q$, $i_\Theta$.
\end{obligation}

\begin{obligation}[Linkage data]
Label and linkage data $L$ which state not merely that the two sides carry the
same names, but that a specific identification is \emph{permitted}.
\end{obligation}

\begin{obligation}[Split indeterminacy budget]
An indeterminacy budget $I$, separated into absorbable renormalization, blurring,
and unabsorbed residual load.
\end{obligation}

\begin{obligation}[Compatible evaluation in the container]
\label{obl:six}
A numerical evaluation $\nu_C$ in the container, \emph{together with} the
compatibility identities the claim needs --- in particular
$\nu_C \circ i_q = \nu_q$ on the stated source domain. These compatibility
clauses are mandatory: the mere existence of a numerical map on $C$ never
discharges this obligation, and a stipulated invariance does not count as a
proof of a compatibility.
\end{obligation}

\begin{obligation}[Order compatibility]
\label{obl:seven}
The implication itself: from the permitted linkage it \emph{follows} that
$\nu_q(\mathrm{Obj}_q) \in H_\Theta$.
\end{obligation}

The load of the certificate sits almost entirely in the last two, and the
distinction between them is the single most useful thing the instrument supplies.
Obligation~6 asks whether the two sides are put on one scale by a stated
compatibility --- not merely housed in one container. Obligation~7 asks whether
the order relation between the two evaluations is derived. The two can fail
separately and for different reasons, and holding them apart is what turns a
diffuse suspicion into two nameable missing statements. Without obligations~6
and~7, cohabitation in a common container is qualitative proximity, not the
numerical statement that the argument needs.

\paragraph{Provenance of the sixth obligation's wording.} Version 0.1 of this
paper stated obligation~6 disjunctively: a numerical evaluation in the container
\emph{or} a provable compatibility. The conjunctive form above is taken from the
methods paper of the series~\cite{GeigerBCM2026}, where the corresponding
contract component was hardened after this paper first appeared. The change is
not cosmetic --- it is the reason two verdicts of version 0.1 are recoded here
(``Changes in version 0.2'') --- and it resolves a tension this paper carried
from the start: the working lemma below already demanded a
``label-compatible, scale-preserving'' evaluation contract, which is more than
the disjunctive form of obligation~6 required.

\paragraph{When an obligation counts as discharged.} Because the obligations differ
in kind, the evidence that discharges them differs too, and saying so in advance is
what keeps the coding symmetric rather than convenient. Obligations~1--5 may be
discharged by an explicit stipulation: a definition, a declared property, a fixed
convention. Obligation~6 requires more than a shared symbol and more than a
shared container: an inspectable assignment of both sides to one evaluation
\emph{and} the compatibility identities that tie that evaluation to the readouts
it is supposed to represent. That both sides are real numbers is not one of them,
and neither is a stipulated invariance. Obligation~7 is recorded as discharged
only where the text
contains a separately identifiable derivation, a named lemma, or an explicitly
delegated result that yields the order relation from the permitted identification.

The negative form of that last rule is the one that constrains this paper most, and
it is deliberately weak: failure to locate such an item, under the source range and
the search terms reported with the finding, is a \emph{documentary non-finding}. It
is not a demonstration that no derivation exists, and it is not a judgment that the
declared properties would be insufficient if the derivation were written out. That
second question is mathematical, and Section~\ref{sec:cannot} states that the
instrument does not decide it. Table~\ref{tab:obligations_taxonomy} provides a synoptic taxonomy of the seven formal obligations and the triggered extension, including their mathematical objects, operational requirements, corresponding ledger rows, and primary failure modes.

\begin{table}[htbp]
\centering
\small
\setlength{\tabcolsep}{3.5pt}
\renewcommand{\arraystretch}{1.15}
\begin{tabular}{@{}lP{0.22\textwidth}P{0.25\textwidth}lP{0.22\textwidth}@{}}
\toprule
Obligation & Mathematical Object & Operational Requirement & Row & Primary Failure Risk \\
\midrule
1 ($q$-readout) & $(\mathrm{Obj}_q, \nu_q)$ & Explicit log-volumetric readout & R0, R2 & Scale loss / implicit typing \\
2 ($\Theta$-target) & $(\mathrm{Obj}_\Theta, H_\Theta)$ & Closed target half-line $H_\Theta \subset \mathbb{R}$ & R0, R3 & Moving target / post-hoc shift \\
3 (Container) & $(C, i_q, i_\Theta)$ & Shared container with typed embeddings & R1 & Missing cohabitation \\
4 (Linkage) & $L \subset \mathrm{Obj}_q \times \mathrm{Obj}_\Theta$ & Permitted identification (protocol) & R4 & Unpermitted label confusion \\
5 (Budget) & $I = (I_{\mathrm{abs}}, I_{\mathrm{res}})$ & Split into absorbed and unabsorbed load & R5 & Hidden second mode \\
6 (Evaluation) & $(\nu_C, \nu_C \circ i_q = \nu_q)$ & Container metric with mandatory compatibility & R6a & Unscaled container / missing identity \\
7 (Order) & $\nu_q(\mathrm{Obj}_q) \in H_\Theta$ & Derived inequality implication & R6b & Asserted membership / non-derivation \\
\midrule
Ext.~E8 & Mode branch $M_2$ & Triggered secondary evaluation channel & E8 & Undisclosed branch selection \\
\bottomrule
\end{tabular}
\caption{Taxonomy of the seven bridge certificate obligations, formal signatures, and audit criteria.}
\label{tab:obligations_taxonomy}
\end{table}

\subsection{The ledger}
\label{sec:ledger}

The obligations are operationalized as nine rows (Table~\ref{tab:bc_ledger_rows}). Each row is a question to be
answered against the wording of a primary text; each carries a recorded
expectation, fixed \emph{before} any text was examined, and a recorded risk --- the
failure mode that the row exists to catch. Recording the expectation in advance is
not decoration: it is what allows a later reader to see which findings confirmed a
prior belief and which contradicted one. The dating is part of the claim and is
therefore stated: the obligations, the rows and their expectations were written down
in May~2026; the four applications reported below were carried out in August~2026.

\begin{table}[htbp]
\centering
\small
\renewcommand{\arraystretch}{1.12}
\begin{tabular}{@{}lP{0.40\textwidth}P{0.17\textwidth}P{0.28\textwidth}@{}}
\toprule
Row & Test & Prior expectation & Risk the row catches \\
\midrule
R0 & Are the two lines separately typed? & yes & apparent equality of distinct things \\
R1 & Is a common container present? & yes & a container does not supply a scale \\
R2 & Is the numerical readout preserved on the $q$-side? & open & the readout loses its object identity \\
R3 & Is the target half-line fixed on the $\Theta$-side? & yes & the half-line is moved after the fact \\
R4 & Is the $j^2$-load addressed? & open & unrecorded $O(\ell^2)$ residual load \\
R5 & Is the indeterminacy budget separated? & partly & blurring hides a second mode \\
R6 & Is order preservation proved? & open & qualitative linkage is not enough \\
R7 & Is approximation weaker than membership? & probably not, for a closed half-line & an apparent escape route \\
R8 & Is a second mode visible or excluded? & open & the core of the question \\
\bottomrule
\end{tabular}
\caption{BC-Ledger v0: Nine audit rows with operational questions, prior expectations, and targeted failure risks.}
\label{tab:bc_ledger_rows}
\end{table}

Three remarks on how to read this table. First, ``open'' in the expectation
column means that no expectation was fixed, not that the answer was expected to
be negative. Second, the rows are not of equal weight. As written, R6 names only
the order question, but it in fact carries obligations~6 and~7 at once --- a
conflation which the split of Section~\ref{sec:splits} undoes, and which is worth
naming, because while a row exists for the order question and none for the
compatibility question, the second is assessed in running prose rather than at
the ledger. R6 is the row on which the instrument was designed to be informative;
R0, R1 and R3 are hygiene rows whose main function is to make a negative finding
in R6 interpretable rather than generic.

Third, the last row is not a core row, and reading it as one distorts the
denominator. R8 asks after a second evaluation mode, which is a feature of the
domain model rather than of the contract: it maps to no obligation; it is active
only where the claim's domain actually offers several modes, branches or
normalizations; and where that trigger is absent it is out of scope rather than
negative. The successor instrument therefore carries it outside the obligation
crosswalk, as the triggered extension~E8, and reports a finding at that label and
never at a contract component~\cite{GeigerBCM2026}. Nothing about the coding
recorded for it below changes; what changes is its status, and that the eight
core rows R0--R7 are the rows a completeness statement is about.

\paragraph{Version label and row namespace.} The ledger stated here is the
certificate ledger fixed in May~2026. Its rows are unprefixed, which was harmless
while it was the only ledger of the series and is not harmless now: the successor
instrument~\cite{GeigerBCM2026} scopes its rows by a version prefix (v1-T for the
transport ledger, v1-F for the full one), and identical numeric suffixes across
instruments do \emph{not} imply identical semantics --- this paper's R8 is the
second-mode row, whereas the row carrying that number in the transport instrument
asks a different question. Every unprefixed row label in what follows is a label
of the ledger of this paper, and the correspondence with the successor rows is a
crosswalk, not an identity. Where a version label is needed, this version assigns
the ledger the name \textbf{BC-Ledger~v0}; the label is introduced here and is
not taken from the series register, which does not yet carry one for it.

\subsection{Refinements produced by use}
\label{sec:splits}

Three of the rows turned out to conflate distinct questions, and the instrument
was sharpened accordingly. The first two refinements are \emph{results of the
applications} and were not part of the certificate as fixed; the third follows
from the hardened reading of obligation~6 stated above. None of the three is a
mere recommendation any longer: in the successor instrument all three are
normative~\cite{GeigerBCM2026}. The case reports in
Sections~\ref{sec:caseA}--\ref{sec:caseD} stay with the original rows wherever
the split makes no difference to the finding, and as a rule report at the split
where it does --- Case~B at R4a and R4b, and Case~A at R6a and R6b. Case~C's findings
likewise separate at the two halves; its report retains the merged row of
version 0.1, with the separation carried in the running text of
Section~\ref{sec:caseC}.

\paragraph{R2a / R2b.} Row R2 mixes two transports. \emph{R2a} asks whether the
readout survives transport \emph{within} one container --- along the link that the
construction itself provides. \emph{R2b} asks whether it survives the passage
\emph{between} containers. Only R2b is the bridge question proper. The distinction
was forced by the comparative discrimination case of
Section~\ref{sec:caseC}, where the two
transports are visibly different in kind, and confirmed by
Section~\ref{sec:caseA}, where R2b collapses into R6 because the identification of
the containers \emph{is} the bridge step.

\paragraph{R4a / R4b.} Row R4 mixes the place where a load is quantified with the
place where it is booked against the bridge. \emph{R4a} asks whether the load is
quantified in the evaluation of the $\Theta$-side at all. \emph{R4b} asks whether
it is booked at the bridge itself. The distinction was forced by
Section~\ref{sec:caseB}: the answer to R4a is yes and the answer to R4b is no, and
a single word for both hides two different findings.

\paragraph{R6a / R6b.} Row R6 mixes the two load-bearing obligations under a
question that names only one of them. \emph{R6a} asks after obligation~6: is the
common evaluation accompanied by the compatibility identities that tie it to the
two readouts? \emph{R6b} asks after obligation~7: is the order relation derived
from the permitted identification? This split was not produced by the
applications but by the hardened reading of obligation~6
(Section~\ref{sec:obligations}). Once a common evaluation no longer discharges
the sixth obligation by itself, a single row cannot carry both verdicts ---
Section~\ref{sec:caseC} returns different answers at the two halves, and
Section~\ref{sec:caseA} returns the same answer for two different reasons;
in either configuration a single merged row would have hidden the structure.

A further refinement is recorded as an open item rather than adopted: obligation~3
(container typing) currently enters the ledger only through R1, which asks whether
a container is \emph{present}. Both Sections~\ref{sec:caseA} and~\ref{sec:caseC}
suggest that the interesting question is whether the typing is \emph{sustained
through the conclusion}, which is a different question and deserves its own row.

\subsection{Verdicts and confidence}

A row receives one of three verdicts, and the vocabulary is fixed because loose
verdict language is how form-level findings turn into truth-level claims (see Table~\ref{tab:verdict_taxonomy}).

\begin{table}[htbp]
\centering
\small
\renewcommand{\arraystretch}{1.12}
\begin{tabular}{@{}lP{0.72\textwidth}@{}}
\toprule
Verdict & Meaning \\
\midrule
\textsc{delivers} & The obligation is discharged by the wording examined. \\
\textsc{partly} & Part of the obligation is discharged and a named part is not. The
  split must be stated, not averaged. \\
\textsc{does not deliver} & The obligation is not discharged by the wording
  examined. This is \emph{not} a claim that the statement is false, nor that it
  cannot be discharged elsewhere. \\
\bottomrule
\end{tabular}
\caption{Verdict vocabulary for certificate evaluations.}
\label{tab:verdict_taxonomy}
\end{table}

Each verdict carries a confidence, and the confidence is about the \emph{textual}
finding, not about the mathematics: high confidence means that the quotations are
unambiguous and that no verdict rests on a doubtful reading. A verdict may be
high-confidence and narrow in reach at the same time; where that occurs, the reach
restriction is stated next to the verdict rather than inside it.

These labels are author-assessed judgments about the textual evidence --- the
clarity of the quotation, the specificity of the anchor, the coverage of the
search, and the reconstruction burden carried. They are not calibrated
probabilities, and no independent second coding was carried out, so no inter-rater
agreement is reported. Read them as a statement of how firmly the wording supports
the entry, not as a measured reliability.

\subsection{Gap typology}
\label{sec:typology}

Gaps are typed, not graded. Grading (``serious'', ``minor'') smuggles in an
assessment the instrument is not entitled to make; typing states what kind of
thing is missing, which is both more useful and more defensible (Table~\ref{tab:gap_taxonomy}).

\begin{table}[htbp]
\centering
\small
\renewcommand{\arraystretch}{1.12}
\begin{tabular}{@{}lP{0.76\textwidth}@{}}
\toprule
Type & Description \\
\midrule
G & A missing or non-reconstructible proof step. \\
Z & A citation that does not carry the statement it is invoked for --- an
  \emph{import} failure. \\
D & Definition drift: a term is used differently from the way it was defined. \\
N & Not an error, but a proof-relevant imprecision at the level of convention or
  notation. \\
\bottomrule
\end{tabular}
\caption{Typology of identified gaps and textual deficits.}
\label{tab:gap_taxonomy}
\end{table}

Types combine. A step whose repair would require both a missing argument and a
resolution of a definitional ambiguity is recorded as G+D, and both parts must be
named.

\subsection{Import-ledger discipline}
\label{sec:importledger}

Every appeal to standard literature is registered as an \emph{import}, with the
exact locus in the published source. The rule attached to it is deliberately
asymmetric:

\begin{quote}
Imports from peer-reviewed literature are treated as trustworthy. What is
examined is their \emph{fit} --- whether the cited statement is the statement the
argument needs --- not their proof.
\end{quote}

This is the operational content of the Z type. It keeps the audit finite: without
it, every audit degenerates into an audit of the whole of arithmetic geometry. It
also keeps the audit honest in the other direction, because an import whose fit
was not checked must be recorded as unchecked rather than passed over.

\subsection{Rules of evidence}
\label{sec:evidence}

Four rules govern what may be written into a ledger row.

\paragraph{No authority argument.} Neither the examined author's assertion that a
step suffices nor a critic's assertion that it does not is evidence. This applies
in both directions and is not a courtesy: an author's self-assessment is an object
of examination, and a critic's verdict is, for the purposes of the ledger, only
the definition of what a negative control would look like.

\paragraph{Verbatim quotation with anchor.} Every finding is carried by a quoted
sentence with a line anchor into a named extract. A finding that cannot be
quoted is not a finding.

\paragraph{``Searched, not found''.} A claim that something is absent from a text
is reported together with the search terms used and the range searched. This is
the rule that makes negative findings falsifiable: a reader who suspects the
search was too narrow can see exactly how narrow it was and repeat it. Every case
report that rests on an absence claim carries such a block: Cases~A, B and~D below
do. Case~C makes no absence claim --- its findings rest on quoted passages that are
present --- and therefore carries none.

\paragraph{Reconstruction discipline.}
\label{sec:reconstruction}
Text extracted from PDF loses relation symbols, Greek letters, exponents and
indices. Where a load-bearing sentence is affected, three things are required:
the reconstruction must be stated explicitly, its derivation must be given from
other passages of the same document, and a cross-check must be exhibited at an
independent place. Where the naive reading of a damaged line would invert the
meaning, that line must be flagged separately. This rule is part of the
instrument, not an apology for it --- and it matters here, because in one of the
four applications a verdict does depend on such a reconstruction
(Section~\ref{sec:reconstruction-A}).

\subsection{The working lemma}

The whole instrument is condensed in one statement, which is a task description
and not a theorem.

\begin{workinglemma}
For the forensic reading of a step of the shape described in
Section~\ref{sec:instrument}, it does not suffice to exhibit a qualitative
relation between the two objects inside a common container. What is required is
an explicit, label-compatible, scale-preserving and order-compatible numerical
evaluation contract $\Brxi$. Any candidate for a second mode $T_2$ must either be
visible inside that contract or be booked as an inadmissible hidden auxiliary
scale.
\end{workinglemma}

Throughout, $T_2$ denotes \emph{the second evaluation mode} --- a second scale,
shift or intertwining alongside the one the argument advertises --- and nothing
else. Row R8 is the ledger row that asks after it.

The status of this statement is: working lemma, not proved. It is to be treated
as a precise task for primary-text examination and for positive and negative
controls --- which is what the remainder of this paper does with it.

\subsection{What the instrument cannot decide}
\label{sec:cannot}

Stated positively, so that it is not mistaken for modesty.

\begin{enumerate}
\item \textbf{It cannot certify soundness.} All nine rows may deliver and the
  argument may still fail for a reason no row asks about.
\item \textbf{It cannot decide mathematical questions of substance.} Whether a
  pair of declared properties would suffice as a contract lemma if it were
  proved, or whether an identical global normalization of two isomorphic copies
  of one number field counts as a provable compatibility --- these are
  mathematical questions. The ledger records only that the text supplies a
  property pair, or a convention, where it would need a lemma.
\item \textbf{It cannot see outside the documents examined.} A row returned as
  not delivered is a statement about a wording. The obligation may be discharged
  in material not examined; where that possibility is concrete, the case reports
  name it.
\item \textbf{It cannot distinguish an unmet obligation from an unmeetable one.}
  Nothing in the instrument licenses the step from ``not derived here'' to ``not
  derivable''.
\end{enumerate}

\subsection{Applying the certificate to a new argument}
\label{sec:recipe}

The instrument is not tied to the four cases below. The following recipe is a
prospective protocol distilled from those applications. It was not completely
executed in every reported case: in particular, Case A leaves hygiene rows open,
and the hygiene-before-R6 rule was formulated after the applications. The recipe
is offered for a new transfer argument; the historical partial applications do
not retrospectively satisfy all its steps.

\begin{recipe}
\begin{enumerate}
\item \textbf{Locate the step.} Identify the sentence at which a qualitative
  relation becomes a numerical assertion. This is usually not where the
  document's own summary says it is; verify the structure against the text and
  record the correction, because the next reader should not have to redo it.
\item \textbf{Name the two sides and the container.} Fill obligations 1--3.
  Record the typing of the embeddings, and record whether the typing is dropped
  anywhere before the conclusion.
\item \textbf{Fill the hygiene rows} R0, R1, R3. If one of them fails, a later
  finding at R6 may be an artifact of the failure rather than an independent
  result.
\item \textbf{Separate obligation~6 from obligation~7.} Ask first whether the
  container carries a numerical evaluation \emph{and} the compatibility
  identities that tie it to the two readouts --- in particular
  $\nu_C \circ i_q = \nu_q$ on the stated source domain. A common evaluation
  alone never discharges obligation~6, and a stipulated invariance is not a
  proof of a compatibility. Only then ask whether the order relation is derived
  from the permitted identification. Conflating the two questions is the most
  common way to produce a finding that is both unfair and uninformative;
  collapsing the first into the bare existence of a common scale is the most
  common way to produce one that is too generous.
\item \textbf{Search for the contract lemma explicitly}, with recorded search
  terms, over the whole document and over the proof range separately. Record the
  result either way.
\item \textbf{Type any gap found} (Section~\ref{sec:typology}) and enter every
  appeal to external literature into the import ledger
  (Section~\ref{sec:importledger}).
\item \textbf{Write the ``searched, not found'' block} before writing the
  verdict, not after. Its function is to constrain the verdict, which it cannot
  do if it is composed to fit one.
\item \textbf{State the reach separately from the verdict.} ``Does not deliver,
  high confidence, restricted to document $X$'' is a complete finding;
  ``does not deliver'' alone is not.
\end{enumerate}
\end{recipe}

% =============================================================================
\section{Case A: the corollary and its proof}
\label{sec:caseA}
% =============================================================================

The first application is to the passage the instrument was built for: the proof
of Corollary 3.12 in the third volume of inter-universal Teichmüller theory~\cite{Mochizuki2021c}. Rows R2, R4, R6 and R8 were filled --- R6 at both halves
of the split introduced in Section~\ref{sec:splits}, and R8 as the triggered
extension row it is rather than as a core row; R0, R1, R3, R5 and R7
were not part of this application and remain open
(Section~\ref{sec:convergence}).

\paragraph{Relation to Part~1 of the series.} This passage has already been located
in print. An earlier, repair-oriented study of the same program~\cite{GeigerGap2026}
places the gap at the transition between the same substeps and reports the same
wording --- that the required inequality is said to ``follow formally''. The present
application neither repeats nor strengthens that diagnosis, and it claims no
priority over it. What it does is narrower: it recodes the documentary question as
one obligation of a stated certificate, so that the same site can be compared with
three further texts under one unchanged instrument. That two routes of a different
kind arrive at the same address is recorded here as a finding about the address, not
as a second detection.

This is therefore a \emph{partial} ledger application, and it is reported as one.
The recipe of Section~\ref{sec:recipe}, which asks for the hygiene rows R0, R1
and~R3 to be filled before R6 is weighed, was written afterwards, out of the
experience of these applications --- as were the row splits of
Section~\ref{sec:splits}, and for the same reason the case reports are not
retrofitted to it. The consequence is stated rather than absorbed: what follows is
a finding about four rows. The confidence recorded with each of them is confidence
in that row, and no full-certificate verdict on this passage is claimed here.

\subsection{Locating the passage}

The substeps in which the transition occurs stand \emph{inside the proof} of the
corollary, not in the remark that recapitulates it. The structure, recorded here
so that it need not be re-derived (Table~\ref{tab:case_a_text_blocks}):

\begin{table}[htbp]
\centering
\small
\renewcommand{\arraystretch}{1.12}
\begin{tabular}{@{}lP{0.30\textwidth}@{}}
\toprule
Block & Extract lines \\
\midrule
Proposition 3.9 (log-volume for packets and processions) & 6365\,ff. \\
Theorem 3.11 (multiradial algorithms), proof & 8666\,ff., proof 9032 \\
Remark 3.11.1 with the declared properties & 9036--9300 \\
Corollary 3.12, statement & 9851--9884 \\
Corollary 3.12, proof & 9886\,ff. \\
Step (x) (indeterminacy balance) & 10237--10283 \\
\textbf{Substeps (xi-a)--(xi-h)} & \textbf{10285--10496} \\
Remark 3.12.2 including the toy model & 10573\,ff., toy model 10733--10813 \\
\bottomrule
\end{tabular}
\caption{Structural blocks and line anchors of the IUT-III Corollary 3.12 proof extract.}
\label{tab:case_a_text_blocks}
\end{table}

Four properties declared in Remark 3.11.1 recur throughout and are abbreviated as
in the source: the input prime-strip link (IPL), the strip-holomorphic
expressibility condition (SHE), algorithmic parallel transport (APT), and hidden
internal structures (HIS).

\subsection{Reconstruction disclosure}
\label{sec:reconstruction-A}

The extract of this document loses relation symbols, exponents, indices and Greek
letters. Unlike in the other three applications, one verdict here \emph{does}
depend on a reconstruction, so the derivation is given in full, per
Section~\ref{sec:reconstruction}.

\begin{enumerate}
\item The target set is defined at l.~10385: the extract reads
  \texttt{R-|log()| d=ef \{ R | -|log()|\} R}, to be read as
  $\mathbb{R}_{\le -|\log(\Theta)|} := \{\lambda \in \mathbb{R} \mid
  \lambda \le -|\log(\Theta)|\} \subset \mathbb{R}$.
\item The corollary's own statement resolves the direction of the relation at
  l.~9882/9884: ``\dots\ i.e., $C_\Theta \ge -1$ for any real number
  $C_\Theta \in \mathbb{R}$ such that
  $-|\log(\Theta)| \le C_\Theta \cdot |\log(q)|$'', hence
  $-|\log(\Theta)| \ge -|\log(q)|$.
\item Together these make l.~10446/10448 unambiguous:
  $-|\log(q)| \in \mathbb{R}_{\le -|\log(\Theta)|}$, ``hence also the
  inequality'' $-|\log(q)| \le -|\log(\Theta)| \in \mathbb{R}$.
\end{enumerate}

The disputed membership therefore stands as an \emph{explicit statement} at
l.~10446 --- with the relation symbols restored as set out above --- and is not an
interpretive inference of ours.

\paragraph{Verification against the page images.} Because a verdict of this section
depends on the restoration, the two affected lines were checked against the
published pages themselves rather than against the extract alone. Both restorations
are confirmed. On p.~184 the page carries the membership
$-|\log(q)| \in \mathbb{R}_{\le -|\log(\Theta)|}$ and, displayed beneath it,
$-|\log(q)| \le -|\log(\Theta)| \in \mathbb{R}$, followed by ``then follows
formally''; on p.~174 the corollary's statement carries
$-|\log(\Theta)| \ge -|\log(q)|$ together with the sentence quoted above. The
symbols the extract dropped are therefore present in the source, and in the form
the reconstruction assumed. This also corrected one detail of our own transcription:
the source writes $C_\Theta$ throughout, not a bare $C$, and the quotations here now
follow it.

One further line is flagged because the naive reading inverts its sense. At
l.~10750 the extract reads \texttt{in the sense that R -h = -2h R}; the intended
reading is $\mathbb{R} \ni -h \neq -2h \in \mathbb{R}$, the ``$\neq$'' having been
lost. The surrounding sentence secures this: the two assignments are ``mutually
incompatible and cannot be considered simultaneously''
\anchor{IUT-III}{10748--10749}. A reader who restores ``$=$'' here would conclude
the opposite of what the passage says --- and the page image settles it: p.~190
carries ``[in the sense that $\mathbb{R} \ni -h \neq -2h \in \mathbb{R}$]''.

\subsection{R2 --- is the \texorpdfstring{$q$}{q}-side readout preserved?}

\textsc{Delivers}, with high confidence --- and the reason is structurally
interesting: at the decisive step no transport takes place at all.

The $q$-side is declared indeterminacy-free in the statement of the corollary
itself:

\begin{quote}\small
``for the procession-normalized mono-analytic log-volume of the image of a
$q$-pilot object \dots\ in the multiradial representation of Theorem 3.11, (i),
\emph{which we do not regard as subject to the indeterminacies} (Ind1), (Ind2),
(Ind3) described in Theorem 3.11, (i), (ii).''
\hfill\anchor{IUT-III}{9869--9873}
\end{quote}

The immediately preceding paragraph attributes exactly the opposite to the
$\Theta$-side \anchor{IUT-III}{9862--9864}. The asymmetry is thus \emph{declared},
not accidental. The $q$-pilot is moreover the fixed reference point of the
construction --- property (IPL) calls it ``the `coric'/`fixed' $q$-pilot
$\mathcal{F}^{\Vdash\times}$-prime-strip'' \anchor{IUT-III}{9117--9122} --- and
property (APT) states that there is no set-theoretic transport to speak of:

\begin{quote}\small
``This parallel transport mechanism \emph{does not consist of a simple instance of
transport of some set-theoretic region} \dots\ \emph{via some set-theoretic
function}. Rather, it consists of a construction algorithm that is simultaneously
valid/executable/well-defined with respect to the arithmetic holomorphic
structures in the domain and codomain of the $\Theta$-link''
\hfill\anchor{IUT-III}{9212--9219}
\end{quote}

The readout is concrete: $-|\log(q)|$ is ``the log-volume of the region that
arises from the representation of the $q$-pilot object''
\anchor{IUT-III}{10387--10390}, on a scale anchored by the statement that
``global arithmetic degrees of objects of global realified Frobenioids may be
interpreted as log-volumes'' \anchor{IUT-III}{10391--10393}.

Obligation~1 is therefore discharged, and obligation~4 is discharged on the
$q$-side --- but \emph{stipulatively}: the identity of the pilot is protected by
the declarations ``coric''/``fixed'' and ``do not regard as subject to'' rather
than produced by a proved relation.

\subsection{R4 --- is the \texorpdfstring{$j^2$}{j\texttwosuperior}-load addressed?}

\textsc{Partly}, with high confidence. The place of the booking is named; the load
itself is not quantified there, and its quantification is explicitly delegated to
the fourth volume.

The load exists and is named: the relevant monoids are generated by the theta
values $\{q^{j^2}\}_{j=1,\dots,\ell^*}$ \anchor{IUT-III}{654--656}. The booking
site is the procession normalization, and what it books is \emph{labels}:

\begin{quote}\small
``the procession-normalized mono-analytic log-volumes [i.e., where the average is
taken over $j \in \mathbb{F}_\ell^*$] \dots\ are \emph{invariant with respect to
the indeterminacies} (Ind1), (Ind2), and have the effect of \emph{converting the
indeterminacy} (Ind3) \emph{into an inequality [from above]}.''
\hfill\anchor{IUT-III}{10261--10268}
\end{quote}

with the averaging attributed to ``the $\mathbb{F}_\ell^*$-symmetry acting on the
labels of the theta values'' \anchor{IUT-III}{10273--10276}. That is a booking of
symmetry, not of magnitude.

For the corollary the load is then explicitly set aside, and the delegation is
declared openly. Property (HIS) names this very application:

\begin{quote}\small
``if one takes the point of view --- \emph{as will be done in Corollary 3.12
below!} --- that one is only interested in considering the \emph{qualitative
logical aspects/consequences} \dots\ then \dots\ one may \dots\ \emph{forget that
this construction algorithm \dots\ has anything to do with theta functions \dots\
or theta values} $[\text{i.e., the } q^{j^2}]_{j=1,\dots,\ell^*}$!''
\hfill\anchor{IUT-III}{9225--9244}
\end{quote}

and, immediately afterwards:

\begin{quote}\small
``unlike Corollary 3.12 below, which only concerns qualitative logical
aspects/consequences \dots, \emph{the explicit computation to be performed in}
[IUT-IV], \S\,1, \emph{of the log-volumes that occur in the statement of Corollary
3.12 makes essential use of the way in which theta values occur} in the
construction algorithm of Theorem 3.11.'' \hfill\anchor{IUT-III}{9263--9267}
\end{quote}

A search of the whole proof range for the theta values, the procession
normalization and the averaging returns hits \emph{only} in Step (x); in the
substeps themselves the load does not reappear.

One point of fairness belongs to the verdict rather than beside it. The failure
mode this row exists to catch --- a load that vanishes in one coordinate and
returns without being booked --- \emph{does not occur here}. What occurs is a
declared division of labour between a qualitative volume and a quantitative one.
For the certificate the finding remains that within the document examined the
load is not booked at the disputed place; the distinction between ``booked
elsewhere, by declaration'' and ``concealed'' is recorded in the source's favor.

\subsection{R6 --- is order preservation proved?}

\textsc{Does not deliver} at both halves of the row, with high confidence. This
is the core finding of the application, and it must be stated in two parts,
because the two obligations fail for different reasons and the row as originally
worded names only the second of them.

\paragraph{R6a --- obligation~6 is not discharged: a common evaluation, no
compatibility identity.} In substep (xi-d) the two sides are
explicitly brought into one evaluation:

\begin{quote}\small
``it is only by forming [a suitable positive tensor power of] the determinant of
the localizations of arithmetic vector bundles \dots\ and then applying the
[suitably normalized \dots] log-volume to various regions \dots\ that we are able
to obtain \emph{completely comparable objects}, namely,
$\mathbb{R}_{\le -|\log(\Theta)|} := \{\lambda \in \mathbb{R} \mid \lambda \le
-|\log(\Theta)|\} \subset \mathbb{R}$; $-|\log(q)| \in \mathbb{R}$''
\hfill\anchor{IUT-III}{10376--10385}
\end{quote}

Together with the log-volume interpretation of arithmetic degrees quoted above,
a numerical evaluation $\nu_C$ exists in the container, and it is stated
explicitly rather than left to be inferred. That much is recorded in the
source's favor. It is not, however, what obligation~6 asks for. The obligation
asks for the compatibility identity $\nu_C \circ i_q = \nu_q$ across the
identification of the two containers, and that identity is not in the text: the
``searched, not found'' block below records the search, and records that the one
compatibility lemma which does exist --- the log-link compatibility of the
log-volumes --- secures compatibility \emph{within} a vertical column rather than
across the identification. Since a common evaluation never discharges
obligation~6 by itself (Section~\ref{sec:obligations}), R6a does not deliver.

\medskip\noindent
Version 0.1 of this paper coded this half the other way, on the disjunctive
wording of obligation~6 then in force, and stated that ``the gap sits at
obligation~7, not at obligation~6''. The quotations are unchanged and no source
was re-examined; the coding follows the hardened wording. See ``Changes in
version 0.2''.

\paragraph{R6b --- obligation~7 is set as a condition, not derived.} Substep (xi-e)
renders (SHE) operationally as a side condition on the admissible operations:

\begin{quote}\small
```expressibility relative to the arithmetic holomorphic structure in the
1-column' [cf.\ (SHE)] amounts precisely to `\emph{expressibility via operations
that are valid/executable/well-defined even when subject to the condition that
the pilot-object log-volume associated to the input data prime-strip \dots\ be
equal to the fixed value} $-|\log(q)| \in \mathbb{R}$'.''
\hfill\anchor{IUT-III}{10412--10420}
\end{quote}

and substep (xi-f) draws the membership from it:

\begin{quote}\small
``the construction of the subset $\mathbb{R}_{\le -|\log(\Theta)|} \subset
\mathbb{R}$ of possible pilot-object log-volumes of output data \emph{is subject
to the condition} that this construction of output data possibilities
\emph{constitutes, in particular, a construction} [perhaps only up to some sort of
`approximation', as a result of various indeterminacies] \emph{of the
pilot-object log-volume of the input data prime-strip}, namely,
$-|\log(q)| \in \mathbb{R}$. The inclusion
$-|\log(q)| \in \mathbb{R}_{\le -|\log(\Theta)|}$, hence also the inequality
$-|\log(q)| \le -|\log(\Theta)| \in \mathbb{R}$ \dots\ \emph{then follows
formally}.'' \hfill\anchor{IUT-III}{10437--10451}
\end{quote}

Substep (xi-g) names the result as ``two \emph{tautologically equivalent} ways to
compute the log-volume of the $q$-pilot object'' \anchor{IUT-III}{10453--10458},
and the author's own recapitulation repeats the pattern: the $q$-intertwining
``may be thought of as a special case of the $\Theta$-intertwining \dots\
Corollary 3.12 then follows, essentially formally''
\anchor{IUT-III}{10723--10731}.

\paragraph{The toy model writes the carrying assumption out as an assumption.}
Without labels the two assignments contradict one another
\anchor{IUT-III}{10747--10750}; the decisive step is then introduced as a
supposition:

\begin{quote}\small
``\emph{Now suppose that one verifies that one may construct} the `assignment with
indeterminacies' \dots\ as a mathematical structure that is intrinsically
associated to the underlying structure of the assignment $\Xi_q$ \dots\
\emph{even if one forgets the labels} `$q$', `$\Theta$' \dots, i.e., \emph{even if
one identifies} $q\mathbb{R}$, $\Theta\mathbb{R}$, in the usual way, with
$\mathbb{R}$ [cf.\ the properties (IPL), (SHE)]''
\hfill\anchor{IUT-III}{10778--10784}
\end{quote}

from which ``one then concludes formally'' the desired bound
\anchor{IUT-III}{10806--10813}.

\paragraph{The finding, in certificate language.} In the real case the supposition
is met by the conjunction (IPL)\,$\wedge$\,(SHE) --- that is, by \emph{properties
attributed to the output of an algorithm}, not by a compatibility lemma
$\nu_C \circ i_q = \nu_q$. At the same place the container typing required by
obligation~3 is deliberately released: the identification of $q\mathbb{R}$ with
$\Theta\mathbb{R}$ is a premise of the inference, not its result. The absent item
is thus one and the same at both halves of the row, and it is absent twice over:
the identity is what obligation~6 requires in order to put the two readouts on
one scale by proof rather than by cohabitation, and it is what obligation~7 would
need in order to carry the order relation across that scale.

\paragraph{Fairness: ``tautologically'' is meant affirmatively.} The source uses
this word approvingly, not as a concession: it speaks of concluding ``nontrivial,
albeit essentially tautological, consequences \dots\ such as the inequality of
Corollary 3.12'' \anchor{IUT-III}{8122--8125}. For its author the tautologous
character is the achievement, not the defect. Reading the word as a
self-refutation would be an authority argument in reverse, and is not done here.
The ledger asks only whether the compatibility identity is stated and whether the
order relation is \emph{derived}; by the wording, neither is.

\subsection{R8 --- is a second mode visible or excluded?}

\textsc{Partly}, with high confidence. Named and qualitatively priced, but not
carried through the bridge as a second scale; not excluded either.

The double structure is stated explicitly: a prime-strip ``\emph{simultaneously
equipped with two intertwinings}, namely, the $\Theta$-intertwining, \emph{up to
indeterminacies}, and the $q$-intertwining'' \anchor{IUT-III}{8114--8125}. The
price of the construction is declared as well:

\begin{quote}\small
``this `solution' may be implemented \emph{only at the cost of admitting the
`indeterminacy' constituted by the upper semi-compatibility of} (Ind3).''
\hfill\anchor{IUT-III}{10019--10031}
\end{quote}

and the content of that notion is given by the text itself: ``upper
semi-compatibility, i.e., in a word, \emph{one-sided inclusions, as opposed to
precise equalities}'' \anchor{IUT-III}{8489--8490}.

What does not happen is that the second mode is carried through the disputed
step. Before the bridge it is resolved into a direction --- (Ind3) is converted
``into an inequality [from above]'' --- and at the bridge the internal structure
is, by (HIS), expressly forgotten. Visible before the step, suspended within it.

A near-miss is recorded rather than suppressed: a remark whose title
(``The log-shifted nature of the localization homomorphism'') suggests a visibly
carried second scale concerns only the localization homomorphism of an earlier
proposition \anchor{IUT-III}{5551--5553} and does not support the row.

\subsection{Result of Case A}

The resulting findings across the audited rows are summarized in Table~\ref{tab:case_a_findings}.

\begin{table}[htbp]
\centering
\small
\renewcommand{\arraystretch}{1.12}
\begin{tabular}{@{}lP{0.62\textwidth}l@{}}
\toprule
Row & Finding & Confidence \\
\midrule
R2 & \textsc{delivers} --- readout definitionally fixed and indeterminacy-free;
  no transport step & high \\
R4 & \textsc{partly} --- site named, load not quantified, quantification
  delegated & high \\
R6a & \textsc{does not deliver} --- common evaluation present and explicit;
  compatibility identity $\nu_C \circ i_q = \nu_q$ absent & high \\
R6b & \textsc{does not deliver} --- membership set via a property pair; follows
  ``formally'' & high \\
R8 & \textsc{partly} --- named and priced, suspended at the bridge & high \\
\bottomrule
\end{tabular}
\caption{Case A: Summary of findings across audited BC-Ledger rows.}
\label{tab:case_a_findings}
\end{table}

The address of the two unmet obligations is a single substep. There, the
membership ``then follows formally'' from the condition that the construction of
the output possibilities \emph{also} constitutes a construction of the input
readout \anchor{IUT-III}{10437--10451}. What is absent is a compatibility
identity $\nu_C \circ i_q = \nu_q$ across the container identification, which
obligation~6 requires in its own right, and, on top of it, a contract lemma that
\emph{derives} the order relation from the permitted linkage, which obligation~7
requires. In their place stands the property pair, whose discharge the toy model
writes out as a supposition. What the text does deliver at this substep is the
common evaluation itself, explicitly --- and that is not nothing; but a common
evaluation without the identity is exactly the cohabitation the certificate was
built to keep apart from a numerical statement.

\paragraph{Searched, not found.} Over the whole extract and, separately, over the
proof range: a compatibility lemma of the form $\nu_C \circ i_q = \nu_q$ across
the container identification (search terms: \texttt{compatib}, \texttt{log-volume},
\texttt{equal}, \texttt{identif}); a numerical value for the approximation term
(\texttt{approximation}, \texttt{epsilon}, \texttt{bound}, \texttt{estimate},
\texttt{Ind3}); an explicit booking of the $O(\ell^2)$ load at the bridge
(\texttt{theta value}, \texttt{procession-normalized}, \texttt{average}); a
visibly carried second evaluation scale at the bridge (\texttt{log-shift},
\texttt{vertical shift}, \texttt{two intertwinings}). One important qualification,
recorded so that the negative finding is not read as carelessness: a genuine
compatibility lemma \emph{does} exist --- the log-link compatibility of the
log-volumes \anchor{IUT-III}{10268--10272} --- but it lies on the wrong axis. It
secures compatibility \emph{within} a vertical column, not across the
identification of the two evaluation containers --- which is what obligation~6
asks for, and what obligation~7 would then have to build on.

\paragraph{Reach.} High confidence attaches to the textual findings, row by row, in
the author-assessed sense fixed in Section~\ref{sec:instrument}. The reach is
restricted twice over: only the third volume was examined here, and only four of
the nine rows were filled, so what is reported is a finding at four rows and not a
verdict of the whole certificate. The quantitative continuation named by the source
itself is the subject of Case B.

% =============================================================================
\section{Case B: the declared quantification site}
\label{sec:caseB}
% =============================================================================

Case A left one concrete possibility open: the load that is not booked at the
bridge might be booked at the site to which the text explicitly delegates it. The
second application opens that site~\cite{Mochizuki2021d} and asks two strictly
separated questions.

\begin{description}
\item[F1] Is the $j^2$/procession load actually quantified there?
\item[F2] Does that section \emph{derive} the disputed membership, or does it
  \emph{presuppose} the corollary?
\end{description}

Keeping them apart is essential, because a positive answer to F1 is easily
mistaken for progress on F2, and it is not.

\subsection{F1: quantified}

\textsc{Quantified}, with high confidence. The chain is written out in the source
and is here only exhibited, not recomputed.

The theta value enters as the input quantity of the computation: ``$\lambda$'' is
set to be ``$\ordd(-)$'' of the element $q_{v_j}^{j^2}$ at the bad places
\anchor{IUT-IV}{1594--1597}. The local term carries $j^2$ visibly, as
$j^2/(2\ell(j+1))$ \anchor{IUT-IV}{1656--1657}; after weighting, $j^2/(2\ell)$
remains \anchor{IUT-IV}{1669--1670}. The procession averaging then evaluates the
sum of squares explicitly:

\begin{quote}\small
``it remains to compute the average over $j \in \mathbb{F}_\ell^\star$.
\emph{By averaging over} $j \in \{1,\dots,\ell^\star = (\ell-1)/2\}$ \emph{and
applying} (E1), (E2), we obtain the `procession-normalized upper bound'
\dots\ $= \frac{\ell+5}{4}\log(d_K) - \frac{\ell+1}{24}\log(q) + \dots$''
\hfill\anchor{IUT-IV}{1675--1684}
\end{quote}

the identity applied being $\frac{1}{n}\sum_{m=1}^{n} m^2 =
\frac{1}{6}(2n+1)(n+1)$ \anchor{IUT-IV}{1379--1380}. After the local estimate the
coefficient survives into the final bound, where it appears as
$-\frac{1}{6}\bigl(1 - \frac{12}{\ell^2}\bigr)\log(q)$
\anchor{IUT-IV}{1775--1776}, and in the statement of the theorem itself
\anchor{IUT-IV}{1352--1353}. The residual $\ell$-dependence is named and discussed
as ``an inevitable consequence of the fundamental role played \dots\ by the
$\ell$-torsion points'' \anchor{IUT-IV}{1909--1912}.

The chain as displayed leaves no term unaccounted for: the load arises only at the
bad
places, the good non-distinguished places contribute zero, the archimedean step is
computed separately, and the final step sums them.

\subsection{F2: presupposes}

\textsc{Presupposes}, with high confidence, and the strongest evidence is
negative. Three complete searches over the whole document return: zero hits for
the $q$-pilot; zero hits for each of the four declared properties on which the
bridge rests in Case A; zero hits for the intertwinings. All occurrences of
``pilot'' in this document are $\Theta$-pilots. A document that, under the terms
searched, never names the $q$-pilot and never mentions the properties carrying the
bridge supplies no derivation of the bridge. The unmet obligation of Case A lies
outside this section's subject matter, as far as those terms reach.

The framing says the same thing positively. The section is introduced as producing
``\emph{more explicit versions} \dots\ \emph{of the log-volume estimates for
$\Theta$-pilot objects obtained in} [IUT-III], Corollary 3.12''
\anchor{IUT-IV}{456--458}, and its main theorem opens with ``\emph{Suppose that we
are in the situation of} [IUT-III], Corollary 3.12''
\anchor{IUT-IV}{1281}.

So does the proof logic. The section supplies the upper side --- ``we thus obtain
the following \emph{upper bound}'' \anchor{IUT-IV}{1773}, ``substituting back
\dots\ we obtain the following \emph{upper bound for} `$C$'''
\anchor{IUT-IV}{1821} --- and imports the lower side:

\begin{quote}\small
``and hence, \emph{by applying the inequality} `$C_\Theta \ge -1$' \emph{of} [IUT-III],
Corollary 3.12, conclude that \dots'' \hfill\anchor{IUT-IV}{1355--1356}
\end{quote}

That imported inequality \emph{is} the disputed membership in constant form: the
corollary's conclusion states $-|\log(\Theta)| \ge -|\log(q)|$, ``i.e., $C_\Theta \ge -1$
for any real number $C_\Theta$ such that
$-|\log(\Theta)| \le C_\Theta \cdot |\log(q)|$''
\anchor{IUT-III}{9878--9884}. The architecture the text writes down for itself is
therefore: this section bounds $C$ from above, the corollary bounds it from below,
and only both together say anything about $\log(q)$. An upper bound on $C$ alone
says nothing.

Finally, the section does not even handle the third indeterminacy itself. At all
three places where it arises, the treatment is delegated backwards: it ``is taken
into account by the fact that we are considering upper bounds [cf.\ the discussion
of Step (x) of the proof of \dots\ Corollary 3.12]'' \anchor{IUT-IV}{1609--1611}.

\subsection{Consequence for the ledger}

Three consequences, and only the first is a change to a verdict.

\begin{enumerate}
\item \textbf{Row R4 splits.} The load is quantified in the evaluation of the
  $\Theta$-side (R4a: yes, in the combined reading of the two volumes) and is not
  booked at the bridge (R4b: no in the third volume, out of scope in the fourth).
  Without the split, one word hides two findings.
\item \textbf{The delegation is verified.} The concern that motivated the row ---
  a load that disappears in one coordinate and returns unbooked --- does not apply
  to the two volumes read together: the theta value disappears in the third and
  \emph{returns named and computed} in the fourth, at exactly the place the third
  had specified. This is recorded in the source's favor.
\item \textbf{The R6 verdict is unchanged at both halves, and its reach is
  extended.} Case A had stated its verdict under the reservation that only one
  volume was examined. For the site named by that volume, the reservation is now
  discharged: there is neither a compatibility identity nor a contract lemma
  there, and no candidate for either. The site is not contradictory but
  thematically uninvolved.
\end{enumerate}

One further observation belongs here, as an observation about architecture rather
than a judgment: by taking $C_\Theta \ge -1$ as an input, the fourth volume makes
visible how much hangs on the membership. Its entire explicit computation produces
an upper bound which, without the lower bound supplied by the corollary, yields no
statement about $\log(q)$ at all.

\paragraph{Searched, not found.} A derivation of the membership anywhere in this
volume (search terms: \texttt{q-pilot}, \texttt{pilot}, \texttt{holomorphic hull},
\texttt{inclusion}, \texttt{follows formally}, \texttt{tautolog},
\texttt{Step (xi)}): not found; the one hit for the substep lies in a later
section and refers back to the third volume. A contract lemma or any discussion of
the property pair: not found, zero hits. A numerical value for the toy model's
approximation term: not found --- and the tempting identification of that term
with the bound computed here is \emph{our} association, not a finding of the text,
which nowhere makes it.

\paragraph{Import ledger.} The verdict of F2 turns on a single appeal to outside
material, and the discipline of Section~\ref{sec:importledger} requires that it be
entered rather than passed over (Table~\ref{tab:case_b_import_ledger}).

\begin{table}[htbp]
\centering
\small
\renewcommand{\arraystretch}{1.12}
\begin{tabular}{@{}P{0.24\textwidth}P{0.66\textwidth}@{}}
\toprule
Field & Entry \\
\midrule
Importing passage & IUT-IV, \S\,1, in the proof of its main estimate
  \anchor{IUT-IV}{1355--1356} \\
Statement imported & the inequality ``$C_\Theta \ge -1$'' of IUT-III, Corollary 3.12 \\
Statement needed & the same inequality: this section produces an upper
  bound for $C$ and needs a lower bound it does not derive \\
Fit & exact --- imported statement and required statement coincide \\
What was checked & the fit, not the proof, per the asymmetric rule of
  Section~\ref{sec:importledger} \\
\bottomrule
\end{tabular}
\caption{Case B: Import ledger entry for the IUT-IV main estimate.}
\label{tab:case_b_import_ledger}
\end{table}

\noindent
No type~Z failure is recorded here: the import carries exactly the statement it is
invoked for. That is precisely why the entry matters. The statement imported is the
one whose derivation obligation~7 asks after, so the appeal is sound and the site
still cannot discharge the obligation.

% =============================================================================
\section{Case C: a comparative discrimination case in a different theory}
\label{sec:caseC}
% =============================================================================

The third application is not a second test of the same text. It applies the
unchanged instrument to an independent alternative program~\cite{JoshiATSIII,
JoshiATSIV} which reaches the same inequality by different means. Its function is
to answer a question that a single application cannot: does the certificate return
the same finding everywhere --- in which case it is a template, not an instrument
--- or does it discriminate?

The label matters and is fixed here in the terms of the methods
paper~\cite{GeigerBCM2026}: this is a \emph{comparative discrimination case}, not
a control. Version 0.1 called it a comparison control, which claimed more than
the case can carry --- a control would have to be a transfer argument whose
seventh obligation \emph{is} discharged, and Section~\ref{sec:convergence}
already recorded that no such case was among those examined. The function
described below is unchanged; only the name is corrected.

It discriminates. Three of the four rows deliver, the fourth delivers only in a
restricted sense, and the unmet obligation sits at a different address and has a
different shape. \emph{Where} that unmet obligation sits was itself revised by two
later corrections (Section~\ref{sec:caseC-nachtraege}); the comparison at the end
of this section states the revised reading rather than the first coding, and the
revision is reported as part of the case.

\subsection{R2: delivered, by a proved relation}

The transport along the construction's own link preserves the readout, and the
$j^2$ factor is carried visibly: the scaling relation between the valuations is
stated as part of a theorem, $|-|_{K'_{y_{w,j}}} = |-|^{j^2}_{K'_{y_{w,1}}}$ for
$j = 1,\dots,\ell^*$ \anchor{ATS-III}{1544--1560}, and is used in the explicit
computation of the theta values' Bloch--Kato classes
\anchor{ATS-III}{6458--6484}. The evaluation scale is concrete: a log-volume with
a stated normalization \anchor{ATS-III}{6562--6596}. Obligation~1 is discharged;
obligation~4 is discharged here by a \emph{proved valuation relation}, not by a
stipulation.

This is the contrast that forced the R2a/R2b split. A second transport also
occurs, between two copies of the ambient data related by a global Frobenius, and
there the identification is \emph{set by a normalization choice}: ``Choose the
same global normalization \dots\ With this choice \dots\ one has''
--- this in the \emph{statement} of Theorem 6.10.1 of ATS-IV
\anchor{ATS-IV}{3186--3189} --- the proof of the resulting equality being, in that
theorem's proof, the single sentence that it ``is a consequence of the choice of
the same global normalization'' \anchor{ATS-IV}{3262--3266}. Same row, two
different kinds of answer.

\subsection{R4: delivered, by a global constraint}

The load arises only at the semistable places, survives into the numerical
evaluation, and its compensation is booked --- the booking instrument being the
product formula:

\begin{quote}\small
``the non-trivial scaling relationship \dots\ at primes $w \in
V^{\mathrm{odd},ss}$ \emph{forces that suitable (re)normalization of valuations
must be introduced} at primes $w \in V_{L'} - V^{\mathrm{odd},ss}$ so that the
product formula \dots\ hold[s] for each normalized arithmeticoid''
\hfill\anchor{ATS-III}{1600--1609}
\end{quote}

Honesty requires the limitation alongside: the proof of that statement is one
sentence, so the \emph{existence} of a consistent compensating renormalization is
inferred from the constraint rather than constructed. The load is located, not
quantified. The failure mode the row exists to catch does not occur.

\subsection{R6: partly --- proved as an inequality, conventional as a bridge}

The inequality itself is a theorem with a proof --- Theorem 9.11.1 of ATS-III, in
its \S\,9.11 \anchor{ATS-III}{6794--6806} --- and
the evaluation in the container is genuinely numerical, computed as a sum of local
contributions \anchor{ATS-IV}{764--769}. Obligation~6 is \emph{partly}
discharged: the evaluation is numerical, but the compatibility holds only between
the sums and is place-wise disclaimed by the source itself --- as the passages
quoted in the next paragraph state in the source's own words. Version 0.1 coded
this half as discharged ``via its first alternative''; under the hardened
wording of obligation~6 the same quotations yield the partial code, and the
reasoning was already in the text (``Changes in version 0.2'').

The bridge that carries the final conclusion is a different matter. Two
evaluations, in two copies related by the global Frobenius, are joined by a middle
equality whose proof, in the proof of Theorem 6.10.1, reads in full:

\begin{quote}\small
``The middle equality is a tautology using the fact that the number fields
provided by $y_0'$ and $y_0$ are identically normalized.''
\hfill\anchor{ATS-IV}{3272--3273}
\end{quote}

while the same text expressly forbids the place-by-place comparison a reader might
use to check it: ``one cannot claim that $z_p \le z_p'$ holds for all primes $p$
without violating the middle equality'' \anchor{ATS-IV}{816--819}, and the
difference by Frobenius at each prime ``\emph{precludes direct comparison of
various local quantities}'' \anchor{ATS-IV}{3208--3211}. Order compatibility thus
exists only between the sums, never place by place.

\subsection{R8: delivered --- and this is why the rest is legible}

The second mode here is the global Frobenius shift, and it is not merely present
but declared, used asymmetrically, interpreted, and flagged as a correction of the
author's own earlier version: the lower bound is computed in one copy and the
upper bound in the other, ``the upper bound is not independent of this choice''
\anchor{ATS-IV}{733--743}, and a bracketed note records that ``[t]he previous
release of this paper had omitted this Frobenius-shifting''
\anchor{ATS-IV}{755}.

This is the strongest positive result of the comparison, and it has an analytic
payoff rather than a rhetorical one: \emph{because} the second mode is disclosed
rather than absorbed into a normalization, the gap at R6 has an address at all.

\subsection{Two later corrections, pulling in opposite directions}
\label{sec:caseC-nachtraege}

Both must be reported, or the case is misrepresented.

\paragraph{Strengthening.} A source examination of the two earlier volumes of the
same series, published separately as a tool supplement of this
series~\cite{GeigerGNC2026}, settled the definitional question left open above: the
normalized object is unique per point, so ``the same global normalization'' is
\emph{definitional}. On that reading the identification is not load-bearing, and
the middle equality rests instead on the by-construction Frobenius stability of
the collated adelic set --- that is, on precisely the mechanism this application
classified as delivered under R8 --- together with a definitional pinning of the
scale. What remains is not a missing contract lemma but a narrower question about
a measure convention.

\paragraph{Weakening.} A substance examination of the theorem carrying R6 ---
companion work of this series, published as Part~5~\cite{GeigerSubstance2026} --- found that
its core step, a containment declared ``clear from the construction'', is not
carried by the construction cited for it. That is a G+D gap: a missing
argument (a supremum statement is used where a measure statement is needed)
together with a structural ambiguity about where one algebraic structure becomes
another. The repair duty is named; the finding is not an error claim.

The two together are best summarized in the words those two examinations used:
the difference in \emph{form} remains, and the advantage in \emph{substance}
shrinks.

\paragraph{The revised coding.} Taken together, the two corrections do not make
this case better or worse than the first coding recorded; they \emph{move} the
deficit. Rows R2, R4 and R8 are untouched by either correction and continue to
deliver. What changes is the split inside the \textsc{partly} of row R6. Before the
corrections, the delivered part was the inequality as a proved theorem and the
missing part was a bridge joined by a conventional identity. After them, the
identity is definitional and therefore not the missing part, while the missing part
sits inside the theorem itself, at its containment step. The verdict stays
\textsc{partly} because the rule is that a split verdict must state its split
rather than average it --- and it is the split that has changed, not the verdict.

That a later and deeper reading revised the study's own first coding is reported
here rather than smoothed away. An entry that a subsequent examination of the same
sources can correct is an entry that can be checked; that is a property worth
having. It is not a claim that the instrument has been validated.

\subsection{What the comparison shows}

Both arguments fail obligation~7 in the strict sense of the ledger, and they fail
it differently --- which is the whole point of a comparative discrimination case (Table~\ref{tab:case_a_c_comparison}).

\begin{table}[htbp]
\centering
\small
\renewcommand{\arraystretch}{1.12}
\begin{tabular}{@{}lP{0.40\textwidth}P{0.40\textwidth}@{}}
\toprule
Row & Case A & Case C \\
\midrule
R2 & delivers --- no transport occurs (declared) & delivers --- proved valuation
  relation; second transport conventional \\
R4 & partly --- site named, quantification delegated & delivers --- booked via the
  product-formula constraint \\
R6 & does not deliver at both halves --- no compatibility identity, and
  membership ``follows formally'' & partly, revised --- compatibility only
  between the sums, place-wise disclaimed; inequality carried by a theorem,
  deficit inside that theorem's containment step \\
R8 & partly --- named, then suspended & delivers --- disclosed and used
  asymmetrically \\
\bottomrule
\end{tabular}
\caption{Comparative discrimination matrix: BC-Ledger findings in Case A (IUT-III) versus Case C (ATS-IV).}
\label{tab:case_a_c_comparison}
\end{table}

The difference that matters is the \emph{shape} of the unmet obligation, and the
corrections above have changed where that shape sits in Case C without changing
that it is one identifiable step. In Case C the deficit is localized: after the
corrections it sits at a single containment step inside the theorem that carries
the inequality, where before it was read as sitting at the conventional identity
joining the two evaluations. In Case A it is areal: a conjunction of two declared
properties, distributed over a remark and four substeps, discharged by an inference
the text itself calls formal. The contrast is between a deficit at one identifiable
step and a deficit distributed over several. It is not a statement about which is
more serious --- the typology of Section~\ref{sec:typology} types gaps and
deliberately does not grade them.

% =============================================================================
\section{Case D: downstream applications}
\label{sec:caseD}
% =============================================================================

The fourth application asks a different question of two papers that \emph{use} the
disputed inequality~\cite{ZhouI, ZhouII}: do they presuppose the certificate,
circumvent it, or introduce an alternative scale? A finding here cannot detect the
gap --- it can only show how the gap propagates.

\paragraph{A note on the label.} These papers are filed in the underlying material
under a name that does not appear in them; the curves they use are
Frey--Hellegouarch curves, and the author is a single author. The label is an
archival artifact and is not used here.

\subsection{Where the bound enters}

In the first paper the disputed inequality appears as Proposition 1.9, its
formula~(1.7), whose entire proof is one sentence: it ``follows from the
$\mu_6$-version of'' the corollary of Case A \anchor{Zhou-I}{728--732}. The formula
number is used exactly once afterwards, in the proof of the next proposition
\anchor{Zhou-I}{753--757}, and Proposition 1.9 is never revisited.
There is no derivation, no intermediate lemma, and no contract lemma of the kind
obligation~7 requires.

Everything the paper contributes of its own lies downstream, on the $\Theta$-side:
a sharper hull bound \anchor{Zhou-I}{482--496}, its own readout of the
$\Theta$-pilot \anchor{Zhou-I}{677--725}, and its own constants in the target
inequality \anchor{Zhou-I}{742--751}, which the paper itself presents as an
improvement over an earlier explicit estimate \anchor{Zhou-I}{1931--1934}.

The second paper never names the corollary at all. It inherits the bridge through
the first paper, at exactly the bridging proposition
\anchor{Zhou-II}{291--292}, and it goes further: a remark tracks where one
hypothesis is used in the four cited volumes and concludes from that tracking that
their results ``still hold for'' a different class of initial data
\anchor{Zhou-II}{267--280}. The disputed substep is never touched.

\subsection{Silence, and a revealing contrast}

Searches over both documents return zero hits for the known critique, for the
indeterminacies, for the multiradial construction, for the pilots, and for every
term of dispute or correctness we searched for. The positioning is implicit
instead, through the word ``unconditionally'' applied to the consequences
\anchor{Zhou-I}{29--30}.

What makes this more than an absence is the contrast. Both papers do practice
explicit reservations about rigour --- the second states plainly that it omits a
step and ``claim[s] that we can always do the similar things presented above''
\anchor{Zhou-II}{153--156} --- but they apply such reservations only to their own
omissions, never to the imported step.

\subsection{No alternative scale}

The log-volume normalization is quoted from the fourth volume of the source theory
rather than replaced \anchor{Zhou-I}{474--480}. The sharper bounds, the separate
readout and the new constants all sit on the sides of the bridge, not in it. One
question is expressly left undecided: whether an explicit factor appearing in the
imported inequality is a renormalization of the same quantity that the corollary
folds into its notation, or a different quantity, could not be decided within the
sources of that application, and is recorded as a follow-up rather than answered.

\subsection{Result of Case D}

Both papers \emph{presuppose} the certificate. Neither circumvents it and neither
introduces an alternative scale for it. The first contributes one element of
obligation~6 --- its own explicit numerical readout on the $\Theta$-side --- but
supplies no compatibility identity, and it takes obligation~7 over wholesale. The unmet obligation is thus not closed downstream
but inherited unexamined, and in the second paper extended to a class of data the
cited volumes do not treat.

\paragraph{Searched, not found.} Both extracts were searched case-insensitively;
the counts are given as first paper\,/\,second paper. \texttt{Scholze} 0/0;
\texttt{Stix} 0/0; \texttt{indetermin} 0/0; \texttt{multiradial} 0/0;
\texttt{Theta-pilot}, \texttt{q-pilot}, \texttt{theta pilot} 0/0;
\texttt{Frobenioid} 0/0; \texttt{controvers}, \texttt{criticis}, \texttt{criticiz},
\texttt{dispute}, \texttt{debate} 0/0; \texttt{gap in}, \texttt{not verif},
\texttt{unverif}, \texttt{validity}, \texttt{correctness} 0/0. Two search terms do
return hits, and neither is a discussion: the masked \texttt{3.12} returns 1/0, the
single hit being the citation at \anchor{Zhou-I}{732}, and
\texttt{log-theta-lattice} returns 1/1, in both cases the title of a bibliography
entry \anchor{Zhou-I}{2583--2584}, \anchor{Zhou-II}{2254--2255}. The reach of this
block is the reach of that word list: a discussion conducted under other
terminology would not have been caught by it.

\paragraph{Version caveat.} The second paper describes itself, in the first line of
its abstract, as an older version to be updated \anchor{Zhou-II}{10}. Every
finding about it is a finding about that version; a later version could in
particular argue the inheritance step differently.

% =============================================================================
\section{What the four applications jointly show}
\label{sec:convergence}
% =============================================================================

\subsection{Four different roles, not four detections}

It would be convenient, and wrong, to say that the finding has been confirmed four
times. The four applications do four different things, and stating them separately
is both more accurate and more informative (Table~\ref{tab:applications_roles}).

\begin{table}[htbp]
\centering
\small
\renewcommand{\arraystretch}{1.12}
\begin{tabular}{@{}lP{0.20\textwidth}P{0.50\textwidth}@{}}
\toprule
Case & Role & What it establishes \\
\midrule
A & \textbf{Detection} & Obligations~6 and~7 are not discharged at a named
  substep: the compatibility identity is not stated, and the order relation is
  not derived. \\
B & \textbf{Closing off a site} & The site to which quantification is delegated
  is explicit and self-contained at the nodes it names, and cannot discharge
  obligation~7, because it consumes
  the disputed inequality as an input. \\
C & \textbf{Comparative discrimination} & The instrument discriminates: three rows
  deliver, and the unmet obligation sits at a different address, in a different
  theory, with a different shape --- an address revised, in the course of this
  study, by a closer reading of the same sources. \\
D & \textbf{Propagation} & Downstream work inherits the unmet obligation without
  examining it, and in one case extends what is inherited. \\
\bottomrule
\end{tabular}
\caption{Systematic roles and distinct methodological functions of the four case applications.}
\label{tab:applications_roles}
\end{table}

Only A is a detection. B is the opposite of a second detection: it establishes
that a specific candidate site is \emph{not} where the obligation is discharged,
which narrows the search rather than repeating it. C is a discrimination case
whose value lies in returning a \emph{different} answer. D says nothing about
whether the
obligation can be discharged; it says that neither of the two versions examined
does so.

Case B carries a second function worth naming, because it is the only internal
evidence here that the instrument can return something other than a deficit. Under
the same criteria that record an unmet obligation in Case A, the text examined in
Case B supplies an explicit quantified chain and an exactly fitting import at the
nodes that carry it: on those questions the answer comes back positive. That is a
\emph{positive documentary comparison}, not a validation --- Case B is not a site
at which obligation~7 could be discharged, so it cannot test the row on which the
certificate is most exposed to the charge of being built for its answer.

\subsection{The address is stable across independent texts}

To provide a complete comparative overview of the empirical findings, Table~\ref{tab:synoptic_cross_case_matrix} synthesizes all four applications across their primary sources, theoretical domains, primary BC-Ledger findings, identified gap types, and methodological functions in the argumentative chain.

\begin{table}[htbp]
\centering
\small
\setlength{\tabcolsep}{3pt}
\renewcommand{\arraystretch}{1.15}
\begin{tabular}{@{}lP{0.17\textwidth}P{0.22\textwidth}P{0.21\textwidth}lP{0.12\textwidth}@{}}
\toprule
Case & Primary Source & Theoretical Context & BC-Ledger Finding & Deficit & Role \\
\midrule
A & IUT-III, Cor.~3.12 & Inter-universal Teichmüller theory & R6a/R6b \textsc{does not deliver} & G (step xi-e) & Detection \\
B & IUT-IV, \S\,1 & Explicit log-volume quantification & R4 \textsc{partly} / Import exact & Z-free & Site closure \\
C & ATS-IV v2, \S\,6.10, Thm.~6.10.1 & Arithmetic Teichmüller spaces & R6 \textsc{partly} (revised) & G (sum-level) & Discrim. \\
D & Zhou-I/II & Downstream Diophantine apps & R6 \textsc{inherited} & G (imported) & Inheritance \\
\bottomrule
\end{tabular}
\caption{Synoptic cross-case matrix of all four applications: scope, primary findings, identified gap types, and epistemic roles.}
\label{tab:synoptic_cross_case_matrix}
\end{table}

What the four do jointly establish is narrower and, we think, more durable than a
count of confirmations. \emph{Two} textually independent bodies of work --- the
theory of Cases~A and~B and the program of Case~C --- examined with one unchanged
instrument, place the load of the argument at the same passage: the step from a
common container to a derived order relation. Neither completes that step, and
neither completes the obligation immediately before it either --- though they
fail there differently, which is what makes the comparison informative. A third
body, Case~D, places nothing: it imports the step from the first, as its own text
says, and can therefore corroborate nothing --- what it adds is the observation
that the obligation travels unexamined. Obligation~6 is unmet in Case~A, where
the compatibility identity is absent outright; partly met in Case~C, where the
compatibility holds between the sums and is place-wise disclaimed; and in Case~D
contributed to on one side only. Obligation~7 is met in none of them in the
strict sense of the ledger.

That the instrument nevertheless separates the cases sharply --- ``does not
deliver'' versus ``partly'' versus ``inherited'' --- tells against reading the
common finding as a mere artifact of the certificate's design. It does not settle
the question. A genuine positive control --- a transfer argument of the same shape
whose seventh obligation \emph{is} discharged --- was not among the cases examined,
and naming that gap is part of the finding rather than an aside to it.

\subsection{An independent formalization project reaches the same site}

During the same period an international formalization project published an interim
report~\cite{LANA2026} which isolates a compatibility problem at the same passage.
Its own words:

\begin{quote}\small
``we isolate a specific compatibility problem at the final stage of the argument:
the relation between the construction arising directly from the $q$-pilot in its
native arithmetic holomorphic structure and the construction obtained by anabelian
and Kummer-theoretic methods through the multiradial procedure. \dots\
\emph{Project LANA has not yet reconstructed a proof of the required
compatibility}.'' \hfill\anchor{LANA}{11--16}
\end{quote}

and, on the status of that statement:

\begin{quote}\small
``the proof of the final numerical inequality \emph{hinges on the compatibility}
(9-1) \emph{which is not manifestly false. However, we \dots\ do not have a proof
of} (9-1) \emph{at this time}.'' \hfill\anchor{LANA}{2988--2990}
\end{quote}

The required statement is an existence claim: two isomorphisms of pointed real
vector spaces with a common target are constructed --- one from the $q$-pilot
``together with the native arithmetic holomorphic structure''
\anchor{LANA}{2807--2812}, the other anabelian and dependent on a choice of
integral structure and ``subject to indeterminacies'' \anchor{LANA}{2783--2803}
--- and what is wished for is ``that \emph{there exists some suitable} $S$
\emph{such that} $\eta^q = \eta_S^{\mathrm{anab}}$'' \anchor{LANA}{2815--2819}.

The report's $R^{\mathrm{ss}}$ quotient identifies adelic measurable sets
when their volumes are equal; this definition does not prove equality of
the two concrete pilot classes or of the maps in (9-1). The equation remains
an open target in the cited report (PDF pp.~46, 48--49). A proposed
object-wise role crosswalk follows (Table~\ref{tab:lana_mapping}); it does
not supply the missing typed implication to obligation~7.

\begin{table}[htbp]
\centering
\small
\renewcommand{\arraystretch}{1.12}
\begin{tabular}{@{}P{0.27\textwidth}P{0.37\textwidth}P{0.23\textwidth}@{}}
\toprule
Obligation & Object in the report & Anchor \\
\midrule
3 --- container, typed embeddings & the common target with its two typed maps &
  2790--2812 \\
4 --- permitted identification & stipulated source identifications plus the choice
  of $S$ & 2937--2939 \\
5 --- split budget & two indeterminacies on the anabelian side, the third as
  ambiguity in the choice of $S$ & 2786, 2973--2974 \\
\textbf{6 --- compatible evaluation} & the common target is by construction the
  volume evaluation quotient, but \textbf{the required equality of the two
  evaluation isomorphisms --- a compatibility identity in the sense of this
  obligation --- is demanded and unproved} & 2790--2794, 2815--2819,
  2988--2990 \\
7 --- order compatibility & no typed implication established here; object,
  normalized evaluation and half-line bound remain to be connected & 2815--2819 \\
\bottomrule
\end{tabular}
\caption{Crosswalk between independent formalization report (Project LANA) objects and certificate obligations.}
\label{tab:lana_mapping}
\end{table}

Three statements about this convergence, and deliberately no fourth.

\begin{enumerate}
\item \textbf{A proposed role correspondence, not validation.} The maps and
  their requested compatibility can motivate a question at obligation~6.
  This is not a formal identification with the certificate contract.
\item \textbf{The typed bridge remains missing.} An implication from (9-1)
  to obligation~7 would require the chosen $S$, the relevant object and
  quotient representatives, a common normalized log-evaluation, and the
  orientation and upper bound of the target half-line to be justified.
  Equality of two abstract maps alone does not supply these premises.
\item \textbf{No order conclusion is derived here.} If those premises and
  an appropriate comparison bound were proved, a compatible evaluation
  could support an order argument. Neither this paper nor the cited interim
  report supplies that complete chain. The correspondence is therefore
  a research task, not a proved one-way implication or a refutation.
\end{enumerate}

The report's own position is also worth recording precisely, because it matches
the verdict vocabulary of Section~\ref{sec:instrument} rather than either
partisan position: the compatibility is ``not manifestly false'', and there is no
proof of it. That is the same distinction the ledger insists on. \emph{Does not
deliver is not refuted} --- and the convergence, such as it is, is a convergence
on where the question sits, not on an answer to it.

\paragraph{What ``independent'' means here, and what it does not.} The word is used
in a strictly documented sense. The certificate and its nine rows were fixed in
May~2026, before any of the four applications reported here were carried out; those
were carried out in August~2026; the interim report appeared in July~2026 and does
not cite this work. That is the whole of the evidence, and it supports exactly one
limited chronological observation --- that the initial ledger predates the
report. This does not rule out influence on later applications or revisions,
and it does not establish the absence of contact or influence in either direction, which is a
claim about people that no reading of the documents could settle, and which is
therefore not made.

Two further limitations belong immediately here. The report is an interim report,
its authors record internal disagreement, and our mapping is an object-wise
assignment, not a proof of equivalence. The mapping also presupposes rather than
exhibits obligations~1 and~2: the two objects and the target half-line are taken as
already typed, and only obligations~3 to~7 are matched in the table above.

\subsection{What is not shown}

\begin{enumerate}
\item \textbf{Five of nine rows are unfilled.} Rows R0, R1, R3, R5 and R7 were not
  part of any of the four applications and are open in all of them. R5 was touched
  only where R4 forced it. A ledger with four rows filled is a partial ledger, and
  the summary tables above should be read as such. The four are also not of one
  status, and the denominator should not obscure it: three of them --- R2, R4 and
  R6, the last reported at both halves of its split --- are core rows, while R8 is
  the triggered extension of Section~\ref{sec:ledger}, which maps to no obligation
  and is out of scope where its trigger is absent. A completeness statement about
  the certificate is a statement about the eight core rows R0--R7, of which three
  were filled.
\item \textbf{No claim about correctness.} Not one proof line was followed in any
  of the four cases, and no estimate was recomputed. In particular Case B
  establishes \emph{that} a computation is carried out and where, not that it is
  right.
\item \textbf{No claim about repairability.} Nothing here says the unmet
  obligation cannot be discharged; one concrete candidate route --- whether some
  text states a theorem deriving the order relation from the declared property
  pair, rather than inferring it formally --- was named as the remaining lever and
  was not pursued.
\item \textbf{No claim about the conjecture.} The word ``gap'' in this paper means
  ``an obligation of a stated certificate that a quoted wording does not
  discharge''. It does not mean that the $abc$ conjecture is open or closed, and
  it does not mean that any theory examined here is wrong.
\end{enumerate}

% =============================================================================
\section{Limitations}
\label{sec:limitations}
% =============================================================================

This section states what the paper does not do, in the terms of the underlying
examinations.

\begin{enumerate}
\item \textbf{No proof was followed.} No theorem of any examined source was
  verified line by line, no estimate recomputed, no constant checked. What was
  classified is the structure of justification.
\item \textbf{Four rows of nine, three of them core rows.} See
  Section~\ref{sec:convergence}. The row splits of Section~\ref{sec:splits} were
  not part of the certificate as fixed --- two arose from use, the third from the
  hardened reading of obligation~6 --- and they are normative only in the
  successor instrument~\cite{GeigerBCM2026}. The case reports therefore stay with
  the original rows wherever the split makes no difference to the finding, so
  that no application is credited with filling a row that did not yet exist when
  it was carried out. Where the split does make a difference they use it and say
  so: Case~B reports at R4a and R4b, and Case~A at R6a and R6b (Case~C's
  separation is carried in its running text against the merged row). Version 0.1
  stated without qualification that the case reports report against the original
  rows; that described this paper's own practice too widely, since Case~B already
  drew its single change of verdict from the R4 split.
\item \textbf{One verdict rests on a reconstruction.} In Case A the load-bearing
  line lost a relation symbol in extraction. The reconstruction and its
  cross-checks are disclosed in Section~\ref{sec:reconstruction-A}; the general
  formula ``no conclusion depends on the placement of a symbol'' would be false
  for this paper and is deliberately not used.
\item \textbf{Case C carries two later corrections in opposite directions}
  (Section~\ref{sec:caseC-nachtraege}), and the case must be read with both. The
  substance examination underlying the second of them is itself a
  form-and-substance examination, not a refutation; it has appeared as
  Part~5~\cite{GeigerSubstance2026}, as has the source examination underlying
  the first~\cite{GeigerGNC2026}.
\item \textbf{Case D describes one version of one document.} The second paper
  declares itself superseded \anchor{Zhou-II}{10}. Its findings are version-bound.
\item \textbf{The convergence of Section~\ref{sec:convergence} is object-wise},
  not formal. No equivalence between the certificate's seventh obligation and the
  report's compatibility statement is claimed or proved, and the implication we do
  state runs in one direction only.
\item \textbf{Mathematical questions of substance are left open by design}, not by
  oversight --- see Section~\ref{sec:cannot}, item~2.
\item \textbf{The instrument is not validated beyond these applications.} Four
  applications in two theories are a small sample. Whether the certificate
  discriminates usefully in other fields is untested; the recipe of
  Section~\ref{sec:recipe} is offered so that it can be tested.
\item \label{lim:biblio} \textbf{Bibliographic discipline.} Every entry was
  checked against its original source before this version (verification pass of
  August 13, 2026). The titles of~\cite{ZhouI} and~\cite{ZhouII} were initially
  omitted --- the material underlying this paper recorded their identifiers,
  author and affiliation but did not quote their titles, and inventing a
  plausible title would be worse than omitting it --- and were subsequently
  added from the arXiv records themselves. The formalization note of April 2026
  that is sometimes mentioned alongside~\cite{LANA2026} is deliberately
  \emph{not} cited: it was secured but not read, and an unread source has no
  place in a bibliography.
\item \textbf{Textual basis are extracts.} All quotations come from
  \texttt{pdftotext} extracts; line references are finding aids, not pagination.
\end{enumerate}

% =============================================================================
\section{Conclusion}
\label{sec:conclusion}
% =============================================================================

The disagreement this paper begins from is not, in the first instance, about
mathematics. It is about what a particular sequence of sentences \emph{is}. For
that kind of disagreement the ordinary instruments are badly suited, because they
ask a question --- is this correct? --- that the two sides answer about different
objects, and because they leave a door open to arguments from authority that
cannot be closed from inside a prose review.

The bridge certificate replaces that question with one a text can answer about
itself: which of seven named obligations does this wording discharge? Around that
question sit disciplines that make the answers auditable rather than merely
confident --- typed rather than graded gaps, an import ledger, recorded search
terms for every negative finding, disclosed reconstructions, and a verdict
vocabulary in which ``not discharged'' cannot quietly become ``refuted''.

Applied four times, the instrument did four different things, and that is the
result worth keeping. It detected two unmet obligations at a named substep and
separated them --- a compatibility identity that is not stated, and an order
relation that is not derived from it --- a distinction that matters, because it
converts a diffuse suspicion into two specific missing statements instead of one
general complaint. It closed off the one site to which the source itself delegates
its quantification, by showing that the computation there is explicit, closed at
the nodes it names, and downstream of the disputed step. It discriminated when
applied to an independent alternative program, returning three delivered rows and
a differently shaped gap at a different address --- and, when a closer reading of the same sources required it,
it revised its own first coding of that case rather than leaving it standing. And
it recorded that two downstream applications inherit the unmet obligation in a
single sentence without examining it.

What would change this picture is a single, statable object: a theorem that
derives the order relation from the permitted identification, rather than
inferring it formally from properties attributed to a construction. That such a
theorem was not found in the material examined is a statement about that material.
Whether it exists elsewhere, or can be proved, is a mathematical question --- and
it is the question this instrument exists to hand over in a form precise enough to
work on.

% =============================================================================
% AI DISCLOSURE -- pipeline standard, included from the canonical template
% (bilingual, carries its own \section*). The template names computational
% scripts and runs; this paper has none, so the restriction paragraph below
% adapts it downwards, which is what the template's own instruction requires.
% Do not drop the restriction: without it the disclosure would claim work not
% done.
% =============================================================================
\phantomsection
\addcontentsline{toc}{section}{AI Disclosure / KI-Offenlegung}
% =============================================================================
% AI DISCLOSURE -- PIPELINE-STANDARD (seit 2026-08-12, User-Freigabe)
% =============================================================================
% Ersetzt die frueheren Stufen L1-L5+ (archiviert unter _archive/).
% EINE Fassung fuer alle Paper dieser Pipeline.
%
% ABWEICHEN NACH UNTEN: Wurde in einem Paper tatsaechlich deutlich weniger
% KI eingesetzt, wird NICHT auf eine andere Datei gewechselt, sondern die
% Formulierung hier abgeschwaecht -- z. B. "in erheblichem Masse in allen
% Phasen" ersetzen durch die tatsaechlich zutreffenden Phasen, und aus der
% Qualitaetssicherungs-Aufzaehlung streichen, was nicht stattfand.
% Der Grundsatz bleibt: ehrlich benennen, was war -- nicht abgrenzen.
%
% Verwendung:
%   \input{../../_templates/disclosure/ai_disclosure_STANDARD}
% =============================================================================

\section*{AI Disclosure / KI-Offenlegung}

% --- German / Deutsch ---

\noindent
KI-Systeme wurden in erheblichem Ma\ss{}e in allen Phasen dieser Arbeit
eingesetzt: Konzeption und Forschungsdiskussion, Literatur- und
Quellenarbeit, mathematische Analyse und Beweisarbeit, Erstellung und
Pr\"ufung der Rechenskripte, Ausf\"uhrung und Auswertung der
Berechnungsl\"aufe, Text, \"Ubersetzung und Satz.

\medskip

\noindent
Qualit\"atssicherung: Einsatz mehrerer unabh\"angiger Modelle
verschiedener Anbieter mit Gegenpr\"ufungen und adversarialen KI-Reviews;
nachvollziehbare, versionierte Rechenskripte mit gespeicherten
Ergebnisdateien und Pr\"ufsummen; fortlaufende Beweis- und
Registerdokumente; strikte Claim-Hygiene (Behauptungen nur mit
dokumentiertem Beleg, Nachtr\"age statt stiller Korrekturen); Changelogs
\"uber alle Versionen.

\bigskip

\noindent\rule{\textwidth}{0.4pt}

\bigskip

% --- English ---

\noindent
AI systems were used extensively in all phases of this work: conception
and research discussion, literature and source work, mathematical analysis
and proof development, creation and checking of the computational scripts,
execution and evaluation of the computational runs, text, translation, and
typesetting.

\medskip

\noindent
Quality assurance: use of several independent models from different
providers with counter-checks and adversarial AI reviews; traceable,
versioned computational scripts with stored result files and checksums;
continuously maintained proof and register documents; strict claim hygiene
(assertions only with documented evidence, addenda instead of silent
corrections); changelogs across all versions.


\medskip

\noindent
\emph{Einschränkung für diese Arbeit:} Zu dieser Arbeit gehören keine
Rechenskripte und keine Berechnungsläufe; die entsprechenden Punkte des
Standardtextes (Erstellung und Prüfung der Rechenskripte, Ausführung und
Auswertung der Berechnungsläufe, versionierte Skripte mit Ergebnisdateien und
Prüfsummen) treffen hier nicht zu. Die Arbeit besteht aus Textbefunden an
lokalen PDF-Extrakten.

\medskip

\noindent
\emph{Restriction for this paper:} this work involved no computational scripts and
no computational runs; the corresponding items of the standard text (creation and
checking of computational scripts, execution and evaluation of computational runs,
versioned scripts with stored result files and checksums) do not apply here. The
work consists of textual findings against local PDF extracts.

% =============================================================================
% BIBLIOGRAPHY
% Provenance of the metadata is stated per entry, in three grades: (i) quoted
% verbatim inside the material underlying this paper; (ii) identifier and
% version verified against the local extract header; (iii) taken from the
% source-checked bibliography of a companion paper of the same program. No
% entry was admitted whose metadata has none of these.
% Mandatory source check completed 2026-08-13 (WebFetch/WebSearch against
% arXiv, Project Euclid/EMS Press, and the cited manuscripts themselves) for
% all 11 entries; see _results/QUELLENCHECK_BIBITEMS_2026-08-13.md. The three
% entries formerly marked provisional (ZhouI, ZhouII, MFHMP) are completed
% below; all other entries were confirmed unchanged.
% =============================================================================
\begin{thebibliography}{99}
\addcontentsline{toc}{section}{References}
\small

% Provenance (i): full bibliographic data quoted verbatim in the underlying
% material, from the bibliography of Zhou I (extract l. 2583-2584).
\bibitem{Mochizuki2021c}
S.~Mochizuki.
\newblock \emph{Inter-universal Teichmüller Theory III: Canonical Splittings of
the Log-theta-lattice}.
\newblock Publ. Res. Inst. Math. Sci. \textbf{57} (2021), no.~1/2, 403--626.
\newblock Cited here as IUT-III; the extract used is of the 2020 institute
version.

% Provenance (ii)+(iii): identifier and role verified in the underlying material;
% journal data from the source-checked bibliography of the companion paper on
% global normalization in the same program.
\bibitem{Mochizuki2021d}
S.~Mochizuki.
\newblock \emph{Inter-universal Teichmüller Theory IV: Log-volume Computations
and Set-theoretic Foundations}.
\newblock Publ. Res. Inst. Math. Sci. \textbf{57} (2021), no.~1/2, 627--723.
\newblock Cited here as IUT-IV; the extract used is of the 2020 institute version.

% Provenance (iii). Volumes I and II are cited only as the context on which the
% examined constructions rest; no finding of this paper rests on them.
\bibitem{Mochizuki2021a}
S.~Mochizuki.
\newblock \emph{Inter-universal Teichmüller Theory I: Construction of Hodge
Theaters}.
\newblock Publ. Res. Inst. Math. Sci. \textbf{57} (2021), no.~1/2, 3--207.

\bibitem{Mochizuki2021b}
S.~Mochizuki.
\newblock \emph{Inter-universal Teichmüller Theory II:
Hodge--Arakelov-theoretic Evaluation}.
\newblock Publ. Res. Inst. Math. Sci. \textbf{57} (2021), no.~1/2, 209--401.

% Provenance (ii) for identifier/version (verified against the local extract
% header), (iii) for the title.
\bibitem{JoshiATSIII}
K.~Joshi.
\newblock \emph{Construction of Arithmetic Teichmuller Spaces III: A `Rosetta
Stone' and a proof of Mochizuki's Corollary 3.12}.
\newblock arXiv:2401.13508, version v4 (2025). Cited here as ATS-III.

\bibitem{JoshiATSIV}
K.~Joshi.
\newblock \emph{Construction of Arithmetic Teichmuller Spaces IV: Proof of the
abc-conjecture}. Preliminary version for comments.
\newblock arXiv:2403.10430, version v2 (2025). Cited here as ATS-IV.

% Provenance (i) for identifier, version, author and affiliation, all quoted in
% the underlying material; titles are NOT quoted there. Titles verified
% 2026-08-13 via WebFetch of the arXiv abstract pages (arxiv.org/abs/2503.14510
% and arxiv.org/abs/2510.05448).
\bibitem{ZhouI}
Z.-P.~Zhou.
\newblock \emph{The inter-universal Teichmüller theory and new Diophantine
results over the rational numbers. I}.
\newblock Institute for Theoretical Sciences, Westlake University.
\newblock arXiv:2503.14510, version v1 (2025). Cited here as Zhou-I.

\bibitem{ZhouII}
Z.-P.~Zhou.
\newblock \emph{The inter-universal Teichmüller theory and new Diophantine
results over rational numbers. II}.
\newblock Institute for Theoretical Sciences, Westlake University.
\newblock arXiv:2510.05448, version v1 (2025). Cited here as Zhou-II. The
document describes itself as an older version to be updated.

% Provenance (i): title and mirror recorded in the underlying material. Date
% corrected 2026-08-13 from "17 July 2026" to "16 July 2026" to match the
% document's own footer, verified via direct WebFetch of the PDF
% (ncatlab.org/nlab/files/LANAProject-Report-July2026.pdf); see
% _results/QUELLENCHECK_BIBITEMS_2026-08-13.md for the discrepancy with
% secondary reports that cite 17 July.
\bibitem{LANA2026}
Project LANA.
\newblock \emph{Interim Report}.
\newblock 16 July 2026 (document date; some secondary reports of the release
cite 17 July 2026). Cited here as LANA.

% Provenance (iii). Cited only as the definition of the negative control; no
% finding of this paper uses it as evidence. Canonical source and date verified
% 2026-08-13 via WebFetch of the manuscript PDF.
\bibitem{ScholzeStix2018}
P.~Scholze and J.~Stix.
\newblock \emph{Why abc is still a conjecture}.
\newblock Manuscript, 16 July 2018. Available at
\url{https://www.math.uni-bonn.de/people/scholze/WhyABCisStillaConjecture.pdf}.

% Provenance (i): title and authors quoted verbatim in the underlying material,
% from the bibliography of Zhou I (extract l. 2589-2591). Publication data
% verified 2026-08-13 via WebFetch of the Project Euclid article page
% (Kodai Mathematical Journal).
\bibitem{MFHMP}
S.~Mochizuki, I.~Fesenko, Y.~Hoshi, A.~Minamide and W.~Porowski.
\newblock \emph{Explicit estimates in inter-universal Teichmüller theory}.
\newblock Kodai Math. J. \textbf{45} (2022), no.~2, 175--236.
\newblock Referred to in Case D as the comparison work for the downstream
constants.

% Own prior work of the same series. Metadata taken from the live Zenodo records
% (checked 2026-08-13); these are the two parts of the series that have appeared.
\bibitem{GeigerGap2026}
L.~Geiger.
\newblock \emph{The Gap in IUTchIII, Corollary 3.12: From Attempted Repair to
Structural Diagnosis}.
\newblock Preprint, version 6.3, Zenodo (2026).
\newblock \url{https://doi.org/10.5281/zenodo.21899747}.
\newblock Part~1 of the \emph{abc} Gap Series; the published forensic study of the
passage examined in Case~A.

% Part 2 of the same series, the methods paper. Published 2026-08-16 while this
% version was being prepared; its hardened sixth contract component is the
% reason for the two recodings declared in "Changes in version 0.2".
\bibitem{GeigerBCM2026}
L.~Geiger.
\newblock \emph{The Bridge-Contract Criterion: A Form Audit for Transport
Arguments, with an Instantiation in the $abc$ Literature}.
\newblock Preprint, version 0.2, Zenodo (2026).
\newblock \url{https://doi.org/10.5281/zenodo.21960477}.
\newblock Concept DOI \url{https://doi.org/10.5281/zenodo.21960476}.
\newblock Part~2 of the \emph{abc} Gap Series; the methods paper, whose hardened
sixth contract component is adopted in Section~\ref{sec:obligations} of this
version.

\bibitem{GeigerSubstance2026}
L.~Geiger.
\newblock \emph{The Joshi ATS Series under the Bridge-Certificate Standard: A
Targeted Substance Audit and an All-Place Rigidity Classification of Effective
Normalization Weights}.
\newblock Preprint, version 0.2, Zenodo (2026).
\newblock \url{https://doi.org/10.5281/zenodo.21956399}.
\newblock Concept DOI \url{https://doi.org/10.5281/zenodo.21956398}.
\newblock Part~5 of the same series; the substance examination cited in
Section~\ref{sec:caseC-nachtraege}.

\bibitem{GeigerGNC2026}
L.~Geiger.
\newblock \emph{Global Normalization in Joshi's Arithmetic Teichmüller Theory:
What the Identification Carries --- and What Carries the Theorem}.
\newblock Preprint, version 1.1, Zenodo (2026).
\newblock \url{https://doi.org/10.5281/zenodo.21924254}.
\newblock Tool supplement of the same series; the source examination underlying the
first of the two corrections in Section~\ref{sec:caseC-nachtraege}.

\end{thebibliography}

\end{document}